\documentclass[anon]{../../journal-templates/NeSy26/nesy2026} % Anonymized submission for OpenReview (double-blind)
% \documentclass[final]{../../journal-templates/NeSy26/nesy2026} % Camera-ready: include author names

\usepackage{amsmath}
\usepackage{amssymb}
\usepackage{mathrsfs,amsopn}
\usepackage{stmaryrd}
\usepackage{multirow}
\usepackage{booktabs}
\usepackage{xcolor}
\usepackage{colortbl}
\let\todo\undefined
\usepackage[hide]{ed}
\renewcommand{\ednote}[1]{\textcolor{red}{[#1]}}
\usepackage{../../quiver}
\usepackage{tikz}
\usetikzlibrary{er,positioning}
\usepackage[most]{tcolorbox}
% Implementation box: \begin{impl}{Title}{Filename.hs} ... \end{impl}
% Plain (no background, no frame, no title -- the title argument is kept for the call
% sites but unused); the code starts at the top, with only the source file name floating
% in the upper right corner.
\newtcolorbox{impl}[2]{%
  blanker, breakable,
  left=0pt, right=0pt, top=0pt, bottom=0pt,
  overlay unbroken and first={%
    \node[anchor=north east, font=\footnotesize\ttfamily, inner sep=1pt, fill=white]
      at (frame.north east) {#2};}
}
% - Haskell source via minted: real, file-backed code (\inputminted) + inline (\mintinline) ---
\usepackage{minted}
\setminted{fontsize=\scriptsize, breaklines=true, breakanywhere=true, tabsize=2}
\newcommand{\hsrc}{../../../NeSyCat.AI.HaskTorch}
\usepackage{../../macros}

% ========================================
% CUSTOM COMMANDS  (paper-specific only; shared macros live in ../../macros.sty)
% ========================================
% bra/ket/braket/ketbra, \oplus, \nesycattorch, \vecTyp: imported from ../../macros.sty
\renewcommand{\P}{\mathrm{P}} % kernel \P (pilcrow) renewed: document-local by design

% ========================================
 % Auto-loaded by the class: amsmath, amssymb, natbib, graphicx, url, algorithm2e
\usepackage{longtable}% for long tables

% arXiv preprint (the [final] build is used for the public arXiv version while under review):
% the paper is NOT yet accepted/published, so suppress the PMLR proceedings banner and the
% (not-yet-)editors. Guarded by \iffinalsubmission, so the [anon] OpenReview build is untouched,
% and the true camera-ready can later restore them by removing this block.
\iffinalsubmission
  \jmlrproceedings{Preprint}{Preprint}%
  \jmlrvolume{}%
  \jmlrworkshop{}%
  \editors{}%
\fi

\title[\nesycattorch]{\nesycattorch: A Differentiable Tensor Implementation of Categorical Semantics for Neurosymbolic Learning}

 % Author block: shown in the [final] build; ignored by the class in [anon] mode (OpenReview stays double-blind).
\clearauthor{\Name{Daniel {Romero Schellhorn}} \Email{daniel.schellhorn@uni-osnabrueck.de}\\
  \Name{Till Mossakowski} \Email{till.mossakowski@uni-osnabrueck.de}\\
  \Name{Bj{\"o}rn Gehrke} \Email{bjoern.gehrke@uni-osnabrueck.de}\\
  \addr University of Osnabr{\"u}ck, Osnabr{\"u}ck, Germany} 

\begin{document}

\maketitle

\begin{abstract}
Neurosymbolic semantics is fragmented: classical, fuzzy, probabilistic and neural systems each define truth by their own inductive rules. NeSyCat, extending ULLER, subsumes them under a single inductive definition of truth, parametric in a strong monad and an aggregation structure on truth-values. NeSyCat has so far lacked an account of predicates and functions learned by neural networks. We provide \nesycattorch as the missing link and interpret computational symbols via neural networks, implementing the framework in probabilistic programming and tensor-based backends. We use the distribution monad for reference semantics and metric evaluation, and complement it by a monad for numerically stable, differentiable training: the lazy log-tensor monad over the log-semiring. For efficient training in batches, we furthermore employ a batch monad. The axioms \emph{are} the source code: written once in monad-based \doNotation, monadic bind performs marginalisation, lazily pruning unneeded branches. On MNIST addition, our HaskTorch, JAX, and PyTorch implementations outperform LTN and DeepProbLog in speed and accuracy, while achieving nearly the accuracy of DeepStochLog. However, unlike DeepStochLog, we stay in a uniform framework that applies to many first-order NeSy approaches. Namely,
the construction is parametric in the monad; instantiating it with, e.g., the Giry monad extends the approach to continuous probability (working out a neural representation here is left for future work).
\end{abstract}


% ============================================================
\section{Introduction}
\label{sec:intro}

Neurosymbolic (NeSy) AI combines the perceptual strength of neural networks with the structured,
verifiable reasoning of symbolic logic. A recurring obstacle is fragmentation: classical, fuzzy,
and probabilistic NeSy systems each come with their own logical language and semantics, so knowledge
bases and learning objectives rarely transfer between them. ULLER - the Unified Language for Learning
and Reasoning \citep{vankriekenULLER2024} -  endows First-Order-Logic (FOL) syntax with three pairwise-independent
semantics - classical, fuzzy, probabilistic - each carrying its own inductive definition of truth.

A recent line of work \citep{schellhornNeSyCatCategorical2026} reformulates all three
semantics as instances of a single \emph{categorical} framework built on \emph{monads}, Moggi's
construct for computational effects in functional programming
\citep{moggiNotionsComputation1991}. The key observation is that an ULLER computation formula
$x := m(T_1,\dots,T_n)\,(F)$, interpreted as ``run model $m$, then bind its result to $x$, then evaluate $F$'', is exactly monadic \textbf{do}-notation. Fixing a strong monad $\mathcal M$ (the
effect) and an aggregated truth-value space $\Omega$ with connectives and quantifiers yields a
\emph{NeSy framework}; classical, fuzzy, probabilistic, LTN, and
possibilistic semantics all reappear as choices of $\mathcal M$ and $\Omega$, evaluated by
\emph{one} inductive definition of truth.
%\ednote{@DS: more about relation to \citep{schellhornNeSyCatCategorical2026}}

For efficiency reasons, we use \emph{lazy} monads. % returning only an abstract syntax tree.
The lifting operation in the distribution monad computes probabilities using marginalization; a lazy monad ensures that marginalization is only done in cases where it is actually needed.

Besides usual FOL function and predicate symbols, we consider computational function symbols $X\to \mathcal M Y$ and computational predicate symbols $X\to \mathcal M \Omega$. At the deep learning level, we need also to consider two-sided computational function symbols $\mathcal M X\to \mathcal M Y$ and two-sided computational predicate symbols $\mathcal M X\to \mathcal M \Omega$. 

We now recall monads and present the monads used in this paper in Table~\ref{tab:monads-nesycat}:



\section{Monads for Computational Effects}
\label{sec:monads}

\begin{definition}[Monad {\citep[\S 3.1]{kohlSchwaigerMonads2021}}]
A monad is given by a triple \minth{(m, return, >>=)} where \minth{m} is a type constructor mapping a type \minth{a} to a type \minth{m a} of computational effects with values from \minth{a}, \minth{return} embeds values into computation and
\minth{>>=} (bind) is used for composition of computations:
\begin{minted}{haskell}
return :: a -> m a
(>>=)  :: m a -> (a -> m b) -> m b
\end{minted}
Here, $c$ \minth{>>=} $f$ first executes computation $c$ over type \minth{a} and passes its value(s) to a function $f$ delivering a computation over type \minth{b}.
A monad needs to satisfy associativity and unit laws. 
\end{definition}

\noindent
Haskell provides the \minth{do} notation as syntactic sugar for composing monadic maps: with it monadic code almost looks like imperative
code, but under the hood there are only pure maps and monads: \minth{do { x <- y; f}} is syntactic sugar for \minth{y >>= (\x -> do { f })}.



\begin{table}[t]
  \centering\small
  \begin{tabular}{@{}cclp{\dimexpr\textwidth-6.6cm\relax}@{}}
  \toprule
  \textbf{Monad} & \textbf{Code} & \textbf{Effect} & \textbf{Description} \\
  \midrule
  $\mathcal D$ & \minth{Dist} & finite probability &
    Finitely supported probability distributions; the reference semantics and metric readout. \\[3pt]
  $\mathcal T$ & \minth{Tens} & logit weights &
    Finite-support Tensor monad $\mathcal{T}\,m = \mathbb R^m$ (leaves are weight tensors) and
    \minth{(>>=)} is the linear pushforward. \\[3pt]
  $\log\mathcal{T}$ & \minth{LogTens} & stable arithmetic &
    \minth{Tens} in logarithmic coordinates, over the log-semiring
    $(\mathbb R, \mathrm{logsumexp}, +)$: numerically stable and differentiable, the monad used in training.
    \\[3pt]
  $\mathcal B$ & \minth{Batch} & batching &
    Reader monad on the batch index object $\underline B$ (Sec.~\ref{sec:training}), parallel
    processing of a mini-batch in training. \\
  \bottomrule
  \end{tabular}
  \caption{Monads in \nesycattorch, in the format of \citet[Table~1]{kohlSchwaigerMonads2021}.}
  \label{tab:monads-nesycat}
  \end{table}


\paragraph{Three monad layers.} In this work, we will use \emph{three} different monads:
\begin{enumerate}
\item the \emph{probability layer}: the distribution monad $\mathcal D$ \citep{schellhornNeSyCatCategorical2026} is the reference semantics. ${\mathcal M} X$ is the space of all finitely supported probability distributions over $X$. \minth{return} $x$ is the distribution assigning all probability mass to $x$. $f$ \minth{>>=} $\rho = \sum_{x\in X}\!f(x)(y)\cdot\rho(x)$. This corresponds to a two-level random process: first $x$ is drawn from $\rho$, then $y$ is drawn from $f(x)$. This results in a marginal distribution for the joint distribution $\hat\rho(x,y):= f(x)(y) \cdot \rho(x)$.
 Distributions over Booleans can be regarded as single
  probabilities (namely those for $\mathit{true}$). 
\item the \emph{tensor layer}: the tensor monad $\mathcal{T}$ works with the same equations. It implements the distribution monad in a differentiable way, but uses real numbers (logits) instead
  of probabilities. It is connected to $\mathcal D$ by a bridge (see Sect.~\ref{sec:mnist} below). The implementation uses \minth{LogTens}, i.e.\ $\mathcal{T}$
  in \emph{logarithmic coordinates} (over the log-semiring) for stable differentiation.
\item the \emph{batching layer}: the batch monad $\mathcal B$, defined in Sect.~\ref{sec:training} below, allows the parallel
  processing of training samples, evaluated together in one run. It carries no probability and no
  geometry, and, dually to the first layer, it never uses the samples' randomness.
  %: the
  %pointwise-evaluation property (Proposition~\ref{prop:pointwise}) holds for any fixed
  %family of inputs.
\end{enumerate}

\section{Syntax and Semantics}
\label{sec:syntax-semantics}
Using these monads, we extend the categorical semantics of \citet[D1.1.]{johnstoneElephant2002} by adding monad symbols and monadic interpretations.
Refining \citep{schellhornNeSyCatCategorical2026}, \nesycattorch organises syntax and semantics into four layers,
each a pair of a \emph{signature} (only symbols) and an \emph{interpretation} (their meaning):
categorical, logical, domain, and grammatical. The first three declare and interpret
symbols; the fourth generates terms and formulas over them and assigns them their monadic
semantics by one induction.

\subsection{Categorical, logical and domain layers}
\label{sec:layer-categorical-logical}

(i) Categorically, we work in the category of sets and maps, $\Set$; computationally it is the category of
(inhabited) Haskell types and maps, on which every Haskell monad of
Section~\ref{sec:monads} is defined. Parameters are drawn from the category $\Tens$ of tensor spaces, which has as objects the Euclidean spaces $\mathbb R^n$ and as morphisms differentiable maps. Implementation-wise, these are implemented as HaskTorch/PyTorch/JAX tensors, where the neural networks live and backpropagation is performed \citep{fongBackpropFunctor2019}. 

The one choice that genuinely varies is the
\emph{effect}: we keep a single \textbf{monad symbol} $\bigcirc$, which allows for the choice of a monad $\mathcal M := \mathcal I(\bigcirc)$ from Table~\ref{tab:monads-nesycat}. This paper uses two readings: \emph{probabilistic}
($\mathcal M := \mathsf{Dist}$), and \emph{tensorial} ($\mathcal M := \mathcal{T}$).

(ii) Logically, the basic truth values are Booleans, $\Omega := \mathbb B = \{\mathbf{False},\mathbf{True}\}$. The
\textbf{connective symbols} $\Conn$ (e.g.\ $\wedge$, $\vee$, $\neg$, $\rightarrow$) are the ordinary
Boolean operations \citep[D1.1]{johnstoneElephant2002} given by $\mathcal I({*}) \colon \mathbb B^{\ari(*)} \to \mathbb B$.
In the neural setting, though, a predicate does not return a plain Boolean but a \emph{monadic} truth value
in $\mathcal M\,\Omega$, for example a distribution over $\mathbb B$, or a vector of logit weights. Hence we lift each connective
to act on $\mathcal M\,\Omega$: bind its arguments, apply the Boolean operation to the plain truth-values
and then return the result (connective clause of Def.~\ref{def:Semantics}).

The only
logical symbols whose interpretation is genuinely monadic are the \textbf{quantifier
symbols} $\mathrm{Quan}$: a quantifier $Q$ takes its body (a map into monadic truth) and aggregates it to a monadic truth value, indexed by a finite (infinitary are future work) object $D$:
\[
  \mathcal I(Q)_D \colon (D \to \mathcal M\,\Omega) \;\longrightarrow\; \mathcal M\,\Omega,
\]

\begin{definition}[Domain signature $\Sigma$]
A domain signature consists of \textbf{domain symbols} $\mathrm{Dom}$, \textbf{variable symbols} $\Var$ (each over a domain symbol,
$x : S$), and \textbf{function symbols} $\Fun$ and \textbf{relation symbols} $\Rel$, each
partitioned into
$\Fun = \Fun^{\mathrm{Tarski}} \sqcup \Fun^{\mathrm{Kleisli}} \sqcup \Fun^{\mathrm{Nesy}}$ and $\Rel = \Rel^{\mathrm{Tarski}} \sqcup \Rel^{\mathrm{Kleisli}} \sqcup \Rel^{\mathrm{Nesy}}:$
\begin{center}\small
\begin{tabular}{@{}llll@{}}
\toprule
\textbf{Kind} & \textbf{Function symbol} & \textbf{Relation symbol} & \textbf{Implementation} \\
\midrule
Tarski  & $f : S_1,\ldots,S_n \to T$ & $R : S_1,\ldots,S_n \to \tau$ & deterministic \\[2pt]
Kleisli & $f : S_1,\ldots,S_n \to \bigcirc T$ & $R : S_1,\ldots,S_n \to \bigcirc\tau$ & effectful \\[2pt]
Nesy    & $f : \bigcirc S_1,\ldots,\bigcirc S_n \to \bigcirc T$
        & $R : \bigcirc S_1,\ldots,\bigcirc S_n \to \bigcirc\tau$ & neural \\
\bottomrule
\end{tabular}
\end{center}
Here, $\tau$ is the type of truth values.
In order to stay close to FOL, we disallow function and relation symbols taking arguments of type $\tau$ or $\bigcirc\tau$.
\end{definition}

\begin{definition}[Domain interpretation $\mathcal I$] A domain interpretation assigns a set
$\mathcal I(S)$ to each domain symbol $S$, and to each function symbol $f$ a map $\mathcal I(f)$ as in
\begin{center}\small
\begin{tabular}{@{}ll@{}}
\toprule
\textbf{Signature} & \textbf{Interpretation} \\
\midrule
$f : S_1,\ldots,S_n \to T$
  & $\mathcal I(f) \colon \mathcal I(S_1,\ldots,S_n) \to \mathcal I(T)$ \\[3pt]
$f : S_1,\ldots,S_n \to \bigcirc T$
  & $\mathcal I(f) \colon \mathcal I(S_1,\ldots,S_n) \to \mathcal M\,\mathcal I(T)$ \\[3pt]
$f : \bigcirc S_1,\ldots,\bigcirc S_n \to \bigcirc T$
  & $\mathcal I(f) \colon \mathcal M\,\mathcal I(S_1) \times \cdots
     \times \mathcal M\,\mathcal I(S_n) \to \mathcal M\,\mathcal I(T)$ \\
\bottomrule
\end{tabular}
\end{center}
where $\mathcal M := \mathcal I(\bigcirc)$ and
$\mathcal I(S_1,\ldots,S_n) := \mathcal I(S_1) \times \cdots \times \mathcal I(S_n)$. Relations similarly, with $\mathcal I(T)$ replaced by $\Omega$. A Tarski symbol is a plain map, a Kleisli symbol a map into
$\mathcal M\,\mathcal I(T)$, and a Nesy symbol takes monadic carriers as input and produces monadic carriers as output.
\end{definition}

\subsection{Grammar}
\label{sec:layer-grammatical}

We write terms $t$ in the \doNotation  of
Section~\ref{sec:monads}, the same notation the implementation uses. To avoid duplication of \doNotation for terms and formulas, we construe formulas $\phi$ as terms of type $\tau$ and connectives as operations on $\tau$.\footnote{This is standard in higher-order logic \citep{church1940formulation,Andrews86}, and with suitable restrictions, it can be used also for first-order logic.} Tarski terms are terms of type $S$ or formulas of type $\tau$. Monadic terms are terms of type $\bigcirc S$ or $\bigcirc \tau$ (the latter are monadic formulas); such terms can be used as terms $t_1,\ldots,t_n,t$ in the \DO-notation. In order to stay close to FOL, the term $t$ in $\mathbf{return}\ t$ must have a Tarskian, i.e.\ $\bigcirc$-free type.
With
$x \in \Var$, $c,f \in \Fun\cup\Rel\cup\Conn$, $Q \in \mathrm{Quan}$, we define terms as follows:
\[
  \begin{aligned}
    t \;&::=\; x \;\mid\; c \;\mid\; f(t_1,\ldots,t_n) && \text{(standard FOL term/formula)}\\
    &\;\mid\; \mathbf{return}\ t \;\mid\;  \DO\{\, x_1 \leftarrow t_1;\, \ldots;\, x_n \leftarrow t_n;\ t \,\} 
    && \text{(monadic term/formula)}\\
    &\;\mid\; Qx(t) && \text{(quantifier, $t,Qx(t):\bigcirc \tau$)}\\
  \end{aligned}
  \]
  Each term $t$ has free variables $\inn(t)$ and (result) type $\out(t)$.
  Term formation has to be type-correct.
  Formulas $\phi$ are terms with $\out\phi=\tau$.
Standard FOL terms can have type $S$ (or $\tau)$ or $\bigcirc S$ (or $\bigcirc \tau$), depending on whether $f$ has a monadic result. Monadic terms are always of type $\bigcirc S$ (or $\bigcirc \tau$).
\doNotation has to be used for the application of connectives to monadic formulas.  
%Since quantifications only apply to monadic formulas \ednote{why?}, they map terms to terms.
We can recover the original ULLER and NeSyCat syntax for terms $T_i$, formula $F$ with variable $x$, and a neural model $m$ in \nesycattorch as follows:
$$[x:=m(T_1,\dots,T_n)]F \quad \equiv \quad \DO \{\ x\leftarrow m(T_1,\dots,T_n);\; F\ \} $$
\noindent
We define the semantics \emph{pointwise} at a 
valuation $\nu$: for the context $[x_1{:}S_1,\ldots,x_k{:}S_k]$, $\nu$ maps variables to values and can be construed as a tuple in
$\mathcal I(S_1) \times \cdots \times \mathcal I(S_k)$, one component per variable. 

The following semantics uses the \doNotation of Section~\ref{sec:monads}.
Subterm values are \emph{always} monadic and therefore always bound with
$\leftarrow$. A \emph{Tarski} symbol is pure, so its result re-enters the monad through
\minth{return}; a \emph{Kleisli} symbol is already monadic, so it stands as the final
\doExpression without \minth{return}; a \emph{Nesy} symbol takes the monadic values
themselves, so no binds occur at all. A quantifier likewise consumes its body as the map $a \mapsto \sem{\phi}(c,a)$ directly, matching its interpretation type from
Section~\ref{sec:layer-categorical-logical}.

\begin{definition}[Semantics $\sem{\cdot}$]
\label{def:Semantics}
Terms and formulas denote maps
\[
  \sem{t} \colon \mathcal I(\inn(t)) \to \mathcal I(\out(t)),
\]
defined by induction on the grammar with $\vec t:=(t_1, \ldots, t_n)$: 

{\setlength{\abovedisplayskip}{4pt}\setlength{\belowdisplayskip}{4pt}%
 \setlength{\abovedisplayshortskip}{4pt}\setlength{\belowdisplayshortskip}{4pt}
\begin{align*}
  \sem{x}(\nu) \;&=\; \nu(x)
    &\qquad
  \sem{f(\vec t)}(\nu) \;&=\;
    \mathcal I(f)\bigl(\sem{t_1}(\nu),\, \ldots,\, \sem{t_n}(\nu)\bigr)
    \\[2pt]
  \sem{c}(\nu) \;&=\; \mathcal I(c)
    &\qquad
  \sem{\mathbf{return}\ t}(\nu) \;&=\;
    \mathbf{return}\ \sem{t}(\nu)
\end{align*}
\begin{align*}
  \sem{\DO\{\, x_1 \leftarrow t_1;\, \ldots;\, x_n \leftarrow t_n;\ t \,\}}(\nu) \;&=\;
    \DO\{\, a_1 \leftarrow \sem{t_1}(\nu);\,\ldots;\, a_n \leftarrow \sem{t_n}(\nu);\\
  &\hphantom{{}=\;}\qquad \sem{t}(\nu, x_1\mapsto a_1,\ldots x_n\mapsto a_n)
     \,\}
\end{align*}
\begin{align*}
  \sem{Qx(t)}(\nu) \;&=\;
    \mathcal I(Q)_{\mathcal I(S)}\bigl(\lambda a.\,\sem{t}(\nu, x\mapsto a)\bigr) \qquad (x : S)
\end{align*}}
%with atomic formulas $R(\vec t)$ exactly as $f(\vec t)$ for the matching kind of $R$.
\end{definition} 

\noindent
The order of independent binds in these clauses is irrelevant: all monads in this paper are
\emph{commutative} \citep{kockMonadsSymmetric1970} ($\mathsf{Dist}$: the independent joint
distribution; $\mathcal{T}$: the outer product), so commutative
connectives stay commutative under the monadic reading.
Section~\ref{sec:mnist} unfolds these clauses step by step on the running example:

% ============================================================
\section{The Running Example: MNIST Addition}
\label{sec:mnist}

We instantiate the MNIST single-digit addition under \emph{distant supervision}: only the
\emph{sum} of two handwritten digits is observed, never the digits. The axiom is
$\forall (x,y,n){:}S \, \bigl(n \mathrel{\mathsf{=}} \mathsf{digit}(x) \mathrel{\mathsf{+}} \mathsf{digit}(y)\bigr)$.
The domain signature $\Sigma$ has sorts $\mathsf{Image},\mathsf{Digit},\mathsf{Nat}$, one
Kleisli symbol $\mathsf{digit}:\mathsf{Image}\to\mathcal M\,\mathsf{Digit}$, and two Tarski
symbols $\mathsf{+}:\mathsf{Digit}^2\to\mathsf{Nat}$, $\mathsf{=}:\mathsf{Nat}^2\to\tau$. Writing
$\mathrm{dig}_\theta:=\mathcal I_\theta(\mathsf{digit})$, the clauses of
Section~\ref{sec:layer-grammatical} are given as:
\[
  \begin{aligned}
\mathbf{do}\ \bigl\{\, d_1 \leftarrow \mathrm{dig}_\theta(x);\ d_2 \leftarrow \mathrm{dig}_\theta(y);\
       \mathbf{return}\,\bigl(n \mathrel{{=}} (d_1 \mathrel{{+}} d_2)\bigr) \bigr\}.
  \end{aligned}
\]
\noindent
Each $\leftarrow$ performs the law of total
probability in $\mathcal{D}$ and the log-space convolution in $\mathcal{T}$ respectively, so this
program is not pseudo-code: it \emph{is} the semantic value, the monad laws guarantee it
equals the literal inductive unfolding, and the identical text compiles as Haskell.

\paragraph{What ``$+$'' becomes.} The Tarski symbol $\mathsf{+}$ is ordinary addition
$\mathsf{sum}:\underline n\times\underline m\to\underline{n{+}m{-}1}$ in the base, where $\underline k=\{0,\dots,k-1\}$; the monad enters
only by \emph{lifting} it: the functor sends this map to $\mathcal{T}(\mathsf{sum})\colon\mathcal{T}(\underline n\times\underline m)\to\mathcal{T}(\underline{n{+}m{-}1})$. Over $\mathbb R$ the carrier is the finitely supported tensor monad
$\mathcal{T}\,m=\mathbb R^m$ (unit $\eta_m(i)=e_i$), a lawful Haskell
\minth{Monad}.\footnote{The finite-dimensional vector-space construction is, strictly, a
\emph{relative} monad on $\mathbf{Fin}\hookrightarrow\mathbf{Set}$ rather than an
endofunctor \citep{altenkirchMonadsNotEndofunctors2015}: its bind sums only over finite index set. The finitely supported version is an ordinary monad.} This lift is the pushforward, which on $\mathsf{sum}$
marginalises a joint vector onto the sum. The two classifiers return a \emph{pair}, joined by the outer product $a\otimes b$, yielding
exactly the discrete convolution:
\[
  \mathcal{T}(\mathsf{sum})(a\otimes b)(s) = \textstyle\sum_{i+k=s} a_i\,b_k = (a * b)(s),\qquad s\in\{0,\dots,18\},
\]
i.e.\ the unnormalized distribution of the sum of two independent digits, exactly what the axiom scores against
(the DeepProbLog reading \citep{DBLP:journals/ai/ManhaeveDKDR21}).

\paragraph{The first two monad layers.}  The digit symbol is \emph{two-sided},
$\mathrm{dig}_\theta \colon \bigcirc\mathsf{Image} \to \bigcirc\mathsf{Digit}$: an observed image enters
as a certain one-hot encoded tensor. Why encode a tensor that is already a tensor? Because the input lives in
$\bigcirc\mathsf{Image}$, a \emph{distribution} over images, and a single observation is just the certain
case: the point mass at that image (one-hot in $\mathcal{T}$, a delta distribution in $\mathcal D$), exactly how one
datum is represented in an empirical data distribution as empirical state. The encoding sends a delta distribution on an element to the one-hot encoding on the same element. The decoding sends logit tensors via a softmax to a distribution. The
probability reading is \emph{defined} as the composition
$\mathrm{dig}_\theta^{\mathcal D} = \mathrm{dec}\circ\mathrm{dig}_\theta^{\mathcal{T}}\circ\mathrm{enc}$:
\[\begin{tikzcd}[ampersand replacement=\&, sep=large]
  {\mathcal{T}(\mathsf{Image})} \&\& {\mathcal{T}(\mathsf{Digit})} \\
  {\mathcal D(\mathsf{Image})} \&\& {\mathcal D(\mathsf{Digit})}
  \arrow["{\mathrm{dig}_\theta^{\mathcal{T}}\ (\text{CNN})}", from=1-1, to=1-3]
  \arrow["{\mathrm{dig}_\theta^{\mathcal D}}"', from=2-1, to=2-3]
  \arrow["{\mathrm{enc}}", from=2-1, to=1-1]
  \arrow["{\mathrm{dec}=\mathrm{softmax}}", from=1-3, to=2-3]
\end{tikzcd}\]
  
 \noindent 
The \textbf{do}-notation derivation above is no
paraphrase: it is the actual Haskell, written once and polymorphic over the monad $m$. The sorts are
plain types, $\mathtt{+}$ and $\mathtt{==}$ are host functions, and
\minth{digit} is the only monad-dependent symbol and the bind
supplies the marginalisation. In Python, we use \mintp{yield} instead of \minth{<-} utilising the generator syntax as syntactic sugar for the same monadic composition as Haskell's \textbf{do}-notation.

\begin{impl}{The abstract formula, written once (monad-polymorphic)}{}
  \small
  \begin{minted}{python3}
class MNistAddition(Example, DistLogTensBridge):
    def formula(self, m, x: Monad[Image], y: Monad[Image], n: Monad[int]) -> Formula[bool]:
        d1 = yield self.digit(m, x)
        d2 = yield self.digit(m, y)
        s = yield n
        return s == d1 + d2
  \end{minted}
\end{impl}

\noindent
Here, apart from the bind \mintp{s = yield n}, which relates only to batching and is explained in Section~\ref{sec:training}, the formula corresponds to the abstract formula in Section~\ref{sec:layer-grammatical}.
The arithmetic is interpreted identically in both monads - plain integer addition - and the only
per-monad choice is how \minth{digit} is read: in \minth{LogTens} the CNN's raw logits become a
leaf; in \mintp{Dist} that same leaf is decoded to a distribution:

\begin{impl}{Two interpretations of \mintp{digit} (one per monad)}{}
  \small
  \begin{minted}{python3}
class MNistAddition(Example, DistLogTensBridge):
    @monad_method
    def digit(self, img: Monad[Image]) -> Monad[int]: ...

    @digit.instance(LogTens)
    def digit_logtens(self, img: LogTens[Image]) -> LogTens[int]:
        model = self.tensor_interpretation.models[type(self).digit]
        return LogTens.bind(img, lambda x: LogDefer(list(range(10)), x, model))

    @digit.instance(Dist)
    def digit_dist(self, img: Dist[Image]) -> Dist[int]:
        return self.decode(self.digit(LogTens, self.enc_dist(img)))
  \end{minted}
\end{impl}

\noindent
The \minth{LogTens} monad represents the type constructor \minth{Tens a = [a -> Real]}
realized over the log-semiring $(\mathbb R,\mathrm{logsumexp},+)$ with finite support given by a list \minth{[a]} of elements of type \minth{a}:

\begin{impl}{The differentiable log-tensor monad (a free monad)}{}
  \small
  \begin{minted}{python3}
class LogTens[A](Monad[A]):                  class Pure[A](LogTens[A]):
    ...                                          value: A

class Bind[A, B](LogTens[A]):                class LogLeaf[A](LogTens[A]):
    dist: LogTens[B]                             support: list[A]
    func: Callable[[B], LogTens[A]]              log_weights: torch.Tensor
  \end{minted}
\end{impl}

\noindent
\minth{LogTens} is used for training since log space is computationally convenient: products of probabilities become sums and marginalisation becomes
log-sum-exp, the numerically stable evaluation \citep{blanchard2021accurately}; weights
range over all of $\mathbb R$ instead of being squashed into $[0,1]$, where
differentiable fuzzy operators degenerate gradient-wise
\citep{kriekenDifferentiableFuzzyOps2022}; and, since softmax is shift-invariant
($\mathrm{softmax}(z) = \mathrm{softmax}(z + c\,\mathbf 1)$), normalising early would
discard the shift, so it is deferred - score now, normalise once at the boundary.

\section{Training with Monads}
\label{sec:training}

Training minimises the \emph{knowledge loss}, the mean negative log-truth of axioms over data,
\[
  \widehat L(\theta) \;=\; \tfrac{1}{N}\sum_{i=1}^{N} -\log \varphi_\theta(s_i),
\]
where $\varphi_\theta(s)$ is the truth value the axiom's semantics assigns to sample $s$
under weights $\theta$ (empirical risk minimisation; equivalently, maximum likelihood on the logical
evidence). This section answers the \emph{computational} question: why \emph{one} batched evaluation
of the axiom computes $\widehat L$, and in particular why an observation may be \emph{bound} in the
\doNotation. 

The obvious
way to compute $\widehat{L}$ runs the \doNotation program $N$ times, once per sample. The
implementation instead runs it on a whole \emph{batch} of samples at once, on tensors
carrying a batch axis.
The finite-sample average $\widehat L$ lives entirely in the first monad - it
exists because the data is a \emph{finite sample} of the world, not because of batching.

The per-sample formula is written once, polymorphic over the first two monad layers (in
$\mathcal D$ for evaluation, in $\mathcal{T}$ for training). The batching layer
is different: it is never a reading of the formula but is applied \emph{outermost}, the training monad is the composite
$\mathcal B \circ \mathcal{T}$.

\begin{definition}[Batch monad]
Let $\underline B = \{1,\dots,B\}$ index the samples evaluated together in \emph{one} run, a
mini-batch of the observations, with $B = N$ for the full sample. The \emph{batch monad}\footnote{The same repeated, conditionally independent structure is called a \emph{plate} in graphical models \citep{buntine1994operations}; cf.\ \texttt{pyro.plate} \citep{bingham2019pyro} and plated factor graphs \citep{obermeyer2019tensor}.} is the
$B$-fold power, i.e.\ the reader monad on $\underline B$:
\[
  \mathcal B\,X \;=\; X^{\underline B} \;\cong\; (\underline B \to X), \qquad
  \eta_{\mathcal B}(x) = \mathrm{const}\,x, \qquad
  (m \bind f)(b) \;=\; f\bigl(m(b)\bigr)(b).
\]
Its bind is the \emph{diagonal}: the continuation of index $b$ sees only index $b$'s value.
\end{definition}

In order to combine the batch monad with other monads, we need to extend it to a \emph{monad transformer}. Monad transformers transform a given monad to a new one by adding extra structure; in this case, by adding a batch dimension.
\begin{proposition}[Batch transformer]
For \emph{every} strong monad $\mathcal M$ the composite
\[
  \mathcal B \mathcal M\,X \;=\; \bigl(\underline B \to \mathcal M\,X\bigr)
\]
is again a monad (the reader monad transformer at $\underline B$ applied to $\mathcal M$), and there are two
canonical monad morphisms into it:
\[
  \mathrm{lift}_{\mathcal M} \colon \mathcal M\,X \to \mathcal B\mathcal M\,X,\;
  \mathrm{lift}_{\mathcal M} = \mathrm{const}
  \qquad\text{and}\qquad
  \mathrm{lift}_{\mathcal B} \colon \mathcal B\,X \to \mathcal B\mathcal M\,X,\;
  \mathrm{lift}_{\mathcal B}\ m = \eta_{\mathcal M} \circ m.
\]
$\mathrm{lift}_{\mathcal M}$ embeds a batch-constant effect;
$\mathrm{lift}_{\mathcal B}$ embeds a batched but \emph{certain} value.
\end{proposition}

A single run in the resulting composite provably yields all per-sample
values of its batch at once (Proposition~\ref{prop:pointwise}).
An observation is a value of the \emph{pure} batch monad, batched but certain,
$m \colon \underline B \to X$, and enters a formula only through
$\mathrm{lift}_{\mathcal B} = \eta_{\mathcal M} \circ (-)$. In the \doNotation this is the
bind $s \leftarrow m$: the bound $s$ is per-sample data carrying no uncertainty.

The training semantics of this paper is the composite $\mathcal B \circ \mathcal{T}$: a
leaf with log-weight tensor of shape $[B, k]$ is the carrier of $\underline B \to \mathcal{T}\,X$
under the \emph{rectangularity} invariant - all $B$ component measures of the batch share one support of
size $k$.
Under this representation $\mathrm{lift}_{\mathcal B}$ is the one-hot embedding: an observed
value becomes the $[B,k]$ one-hot log-weight leaf, the log-space image of $\eta_{\mathcal M}$
(the numerical floor at $\epsilon$ stands in for $\log 0 = -\infty$).

% ============================================================
\section{Inference and Evaluation}
\label{sec:eval}

The inferential level turns the formula into an optimisation problem: the knowledge loss
$\widehat L$, a data loss $L_{\mathrm{data}}$, and their convex
combination. MNIST addition is
pure distant supervision: Adam minimises only the knowledge loss
$\widehat L(\theta)$.

We run the same specification in two backends, Haskell (HaskTorch) and Python (JAX). In Python JAX the log-space convolution is JIT-compiled and
differentiated. All numbers below come from the JAX backend on a single A100 GPU, averaged over $15$
seeds. We never use digit labels: the network learns to read the digits (about $97\%$ digit accuracy)
from the observed sums alone.

\begin{table}[h]
  \centering
  \caption{Single-digit MNIST addition: \nesycattorch's two variants, $3000$ training pairs. The DPL-style sweep is longer only because its harness uses batch $2$ (${\sim}10\times$ more steps) - the network is not slower, its per-step cost is even a little lower.}
  \label{tab:mnist}
  \small
  \begin{tabular}{@{}lcc@{}}
    \textbf{Metric} & \textbf{LTN-style} (100n. head, ELU) & \textbf{DPL-style} (120n. head, ReLU) \\
    \midrule
    Sum accuracy (test)                  & $94.6 \pm 0.6$\,\% & $94.2 \pm 0.7$\,\% \\
    Sum accuracy (train)                 & $\approx 100.0$\,\% & $\approx 100.0$\,\% \\
    Train / test step (batch $32$) & $0.52$\,/\,$0.20$\,ms & $0.44$\,/\,$0.20$\,ms \\
    Wall-clock, $15$-seed sweep  & $\approx 1.5$\,min & $\approx 5.6$\,min \\
    Training harness                     &$32$ batches, $20$ epochs & $2$ batches, to convergence. \\
  \end{tabular}
  \end{table}

Table~\ref{tab:addition} puts \nesycattorch next to LTN \citep{DBLP:journals/ai/BadreddineGSS22},
DeepProbLog \citep{DBLP:journals/ai/ManhaeveDKDR21}, and for two digits, logLTN
\citep{badreddineLogLTN2023}, NeuPSL and DeepStochLog. For the single-digit case we give each \nesycattorch
variant the same CNN, batch size and number of epochs as the baseline it is compared to. \nesycattorch is ahead in both matchups: $94.6$ vs LTN's $93.5$, and $94.2$
vs DeepProbLog's $92.2$. The real gap is a little larger, because LTN reports only its $10$ best of
$15$ runs (about one in five gets stuck early on), while we average all $15$. Although this issue has been solved by working in log space \citep{badreddineLogLTN2023}, we outperform even logLTN in both accuracy and speed. In general, speed is not even an issue: at batch $32$ our step times (Table~\ref{tab:mnist}) are well under the $5.36$/$3.44$\,ms
LTN reports on an older V100 GPU.

\begin{table}[h]
\centering
\caption{Test sum accuracy (\%), mean\,$\pm$\,std, on MNISTAdd for $N$ digits.  \nesycattorch averages $15$ seeds; LTN averages its $10$
\emph{best} of $15$. ${}^1$ from \citep{badreddineLogLTN2023}, ${}^2$ from \citep{kriekenANeSI2023}, ``T/O'': timeout; ``-'': not reported}
\label{tab:addition}
\small
\begin{tabular}{@{}lrccccc@{}}
                                                          &                         & \textbf{$N{=}1$}        & \multicolumn{2}{c}{\textbf{$N{=}2$}} & \textbf{$N{=}4$}                 \\
  \cmidrule(lr){3-3}\cmidrule(lr){4-5}\cmidrule(lr){6-6}
  \textbf{Method}                                         & \textbf{Trainings:}     & $3$k                    & $1.5$k                               & $15$k                   & $7.5$k \\
  \midrule
  \multicolumn{2}{l}{ \textbf{\nesycattorch (LTN-style)}} & $\mathbf{94.6 \pm 0.6}$ & $\mathbf{89.7 \pm 0.7}$ & $95.7 \pm 0.5$                       & $92.0 \pm 0.6$                   \\
  \multicolumn{2}{l}{\textbf{\nesycattorch (DPL-style)}}  & $94.2 \pm 0.7$    & $89.2 \pm 0.7$    & $95.8 \pm 0.6$                 & $91.8 \pm 0.8$                   \\
  \multicolumn{2}{l}{LTN}                                 & $93.5 \pm 0.3$ ${}^1$   & $88.4 \pm 1.0$ ${}^1$   & $95.4 \pm 0.3$          ${}^1$       & T/O                     ${}^2$         \\
  \multicolumn{2}{l}{logLTN}                              & -                       & $88.3 \pm 0.8$ ${}^1$   & $95.6 \pm 0.5$          ${}^1$       & -                                      \\
  \multicolumn{2}{l}{DeepProbLog}                         & $92.2 \pm 1.6$ ${}^1$   & $87.2 \pm 1.9$ ${}^1$   & $95.2 \pm 1.7$          ${}^2$       & T/O                     ${}^2$         \\
  \multicolumn{2}{l}{DeepStochLog}                        & -                       & -                       & $\mathbf{96.4 \pm 0.1}$ ${}^2$       & $\mathbf{92.7 \pm 0.6}$ ${}^2$         \\
  \multicolumn{2}{l}{A-NeSI}                              & -                       & -                       & $96.0 \pm 0.4$          ${}^2$       & $92.6 \pm 0.8$          ${}^2$        
\end{tabular}
\end{table}

\paragraph{Longer numbers.} The $N{=}4$ column uses two four-digit numbers (sums up to $19{,}998$). \nesycattorch, DeepProbLog and (in the non-recursive and mass-conserving case) also DeepStochLog
optimise the same marginal likelihood, so they end up near that cap. \nesycattorch
still runs at $N{=}4$ ($92.0\%$, just below the cap) where DeepProbLog and LTN time out, with
DeepStochLog and A-NeSI a little ahead.

% ============================================================
\section{Related Work}
\label{sec:related}

\textbf{ULLER} \citep{vankriekenULLER2024} leaves the
semantics split into duplicated inductive definitions. The categorical framework of
\citet{schellhornNeSyCatCategorical2026} instead derives one definition parametric in a monad. But they do not implement neural learning, which we do here with \textbf{\nesycattorch}. \textbf{LTN}
\citep{DBLP:journals/ai/BadreddineGSS22} uses \emph{fuzzy} semantics and thereby avoid probabilistic semantics and hence do not derive any distributions.
Further, \textbf{logLTN} \citep{badreddineLogLTN2023} moves LTN's operators to log space
by composing $\log$ with softmax, whereas here log space is the native carrier and the
softmax is confined to the decode bridge.

\textbf{DeepProbLog} \citep{DBLP:journals/ai/ManhaeveDKDR21} trains on the same exact likelihood we do,
but compiles the program to a Sentential Decision Diagram and scores it by weighted model counting;
\citet{kriekenANeSI2023} show this counting is \#P-hard, so it times out at $N{=}4$
(Table~\ref{tab:addition}). \textbf{DeepStochLog} \citep{wintersDeepStochLog2022}  is the most accurate baseline, marginally ahead of both A-NeSI and \nesycattorch. It
 avoids DeepProbLog's blow-up by writing the sum as a stochastic grammar (a random walk rather than a random graph) which is cheaper but can lose mass on failing derivations. However, DeepStochLog changes the aggregation structure to grammar-derivation probability, while our approach keeps the distribution marginal known from DeepProbLog, yet achieves scalability via laziness and tensor structure. DeepStochLog's grammar-based approach corresponds to formulas dynamically adapted to the input (e.g.\ input length); we can achieve this by using \nesycattorch formulas inside Python or Haskell programs (that e.g.\ adapt the number of \textbf{do} binds to the number of input digits).

 \textbf{A-NeSI} \citep{kriekenANeSI2023}
scales instead by training a network to approximate the counting, which gives up exactness and requires additional factorization preparations, which are example specific and therefore not generalizable. \nesycattorch
needs none of this: the same marginalisation is just the monadic bind and by changing the monad the one framework also gives the classical and fuzzy semantics, more aligned with standard first-order logic.
 The
differentiable reading connects to the logic of differentiable logics
\citep{slusarzDifferentiableLogics2023} and to categorical deep learning
\citep{gavranovicCategoricalDeepLearning2024,fongBackpropFunctor2019}, with broader motivation in the
survey of \citet{smetDeepSeaProbLog2023}. 

% ============================================================
\section{Conclusion}
\label{sec:conclusion}

NeSyCat is a general unifying neurosymbolic framework for reasoning and learning in first-order logic that is parameterized over a computational monad. Monads in NeSyCat have covered different theoretical frameworks, such as discrete and continuous probabilities and nondeterminism. With NeSyCat Torch, we provide the first neural implementation of this general framework, i.e.\ using neural networks as realisations of computational function and predicate symbols. We concentrate on finitely supported distributions here. We provide implementation at Haskell: \url{https://anonymous.4open.science/r/nesycattorch-hs/}, Python JAX: \url{https://anonymous.4open.science/r/nesycattorch-jax/}, and Python PyTorch: \url{https://anonymous.4open.science/r/nesycattorch-py/}.

It turns out that for these implementations, we need the finitely supported distribution monad plus two other monads: for (log-scale) tensors and for training batches. NeSyCat Torch thus scales the generality of NeSyCat to the neural implementation level, while simultaneously obtaining competitive results, in terms of accuracy and training time, for the classical MNIST digit addition example. Our approach also scales to multi-digit addition. This is achieved by using lazy monads that defer evaluation of marginalized probabilities to cases where really needed. Future work will study how well this generalises to more complex examples. 
The generalization of our implementation to other monads like continuous probability and to infinite domains is left to future work as well. Note that an efficient neural representation of continuous probability is non-trivial. 

%\textbf{Categorical
%probability} with the Giry monad \citep{giry1982categorical}, quasi-Borel spaces
%\citep{heunenConvenientCategory2017}, and Markov categories
%\citep{fritz_synthetic_2020} supplies the continuous-probability and infinite-domain story.

\bibliography{../../daniel_zotero, ../../additional_citations}

\newpage
\appendix
\section{Categorical Background}\label{app:cat}

\paragraph{Monads.} Categorically, a monad on a category $\mathcal C$ is a functor $T \colon \mathcal C \to
\mathcal C$ with natural transformations $\eta \colon \mathrm{id}_{\mathcal C} \Rightarrow
T$ (\emph{unit}) and $\mu \colon TT \Rightarrow T$ (\emph{multiplication}) satisfying
$\mu \comp \eta T = \mathrm{id} = \mu \comp T\eta$ and $\mu \comp T\mu = \mu \comp \mu T$. The programming definition above corresponds
one-to-one to the equivalent presentation as a \emph{Kleisli triple} $(T, \eta, (\cdot)^{\mathcal M})$,
where $f^{\mathcal M} \colon TA \to TB$ for $f \colon A \to TB$ satisfies $\eta_A^{\mathcal M} =
\mathrm{id}_{TA}$, $f^{\mathcal M} \comp \eta_A^{\mathcal M} = f$, and $g^{\mathcal M} \comp f^{\mathcal M} = (g^{\mathcal M} \comp f)^{\mathcal M}$: \minth{return} is $\eta$, \minth{(>>=)} is the
Kleisli lift $(\cdot)^{\mathcal M}$ (applied flipped), and the \doNotation of Section~\ref{sec:monads}.
is its syntactic sugar.


  \paragraph{States.} We work over a \emph{concrete} Cartesian category $\mathcal{C}$: objects are sets equipped
  with structure (for example measurable spaces, tensor spaces or also plain sets), morphisms are
  structure-preserving maps, finite products exist, and the terminal object $1$ is the
  one-element set. Effectful maps are always written explicitly as maps
  $f \colon S \to \mathcal M\,T$ of the chosen strong monad $\mathcal M$.
  Because $\mathcal{C}$ is concrete and Cartesian, a \emph{state} on $S$, formally a Kleisli
  point $1 \to \mathcal M\,S$, is the same thing as an \emph{element} of $\mathcal M\,S$; we
  use this identification throughout and simply write $D \in \mathcal M\,S$. 
 
  \paragraph{Lattice-Semirings.}
\begin{definition}
A \textbf{lattice-semiring} is a commutative semiring
$(S, \oplus, \otimes, 0, 1)$ with a lattice order $\le$ on the same
carrier: meet $\wedge$, join $\vee$, bounds $\bot$, $\top$ where they
exist, and $\oplus$, $\otimes$ monotone in each argument.
\end{definition}
\noindent
Each row
carries the lattice ($\bot$, $\top$, $\vee$, $\wedge$) and the pair
($0$, $1$, $\oplus$, $\otimes$) sharing the order:

\begin{center}\small
\label{tab:semirings}
\begin{tabular}{@{}ll|llll|llll@{}}
semiring & carrier & $\bot$ & $\top$ & $\vee$ & $\wedge$ & $0$ & $1$ & $\oplus$ & $\otimes$ \\
\hline
Boolean & $\mathbb B = \{0,1\}$ & $0$ & $1$ & $\mathtt{or}$ & $\mathtt{and}$ & $0$ & $1$ & $\mathtt{or}$ & $\mathtt{and}$ \\[2pt]
probability & $\mathbb P = [0,1]$ & $0$ & $1$ & $\max$ & $\min$ & $0$ & $1$ & $p {+} q {-} p\,q$ & $\cdot$ \\[2pt]
mass & $\mathbb M = \mathbb R_{\ge 0}$ & $0$ & $\infty$ & $\max$ & $\min$ & $0$ & $1$ & $+$ & $\cdot$ \\[2pt]
log & $\mathbb L = \mathbb R \cup \{-\infty\}$ & $-\infty$ & $\infty$ & $\max$ & $\min$ & $-\infty$ & $0$ & $\mathrm{lse}$ & $+$ \\
\end{tabular}
\end{center}

\paragraph{Semiring Monads.} A
semiring $S$ yields a monad $\mathcal M_S$: on a set $X$, a state is a
finitely supported $S$ weighted combination of elements of $X$, and bind
is matrix multiplication over $S$. This is a monad on $\Set$ for every
semiring $S$, with algebras the $S$ semimodules
(\citealp[Lemma~5.1.5]{jacobsIntroductionCoalgebra2016};
\citealp[Ex.~2.2]{reggioCodensity2020};
\citealp{altenkirchMonadsNotEndofunctors2015}); the semiring laws of $S$
become the monad laws: addition marginalises, multiplication chains,
distributivity is bind associativity. This gives
$\mathcal T := \mathcal M_{\mathbb M}$ and
$\log\mathcal T := \mathcal M_{\mathbb L}$.  Since $\log$ is a semiring isomorphism
$(\mathbb R_{\ge 0}, +, \cdot) \cong
(\mathbb R \cup \{-\infty\}, \mathrm{lse}, +)$
\citep{mohriWeightedAutomata2009}, we have
$\mathcal T \cong \log\mathcal T$. The distribution monad is the
mass-one submonad $\mathcal D := \{\rho \in \mathcal M_{\mathbb M} \mid
\Sigma\rho = 1\}$ which is \emph{not} $\mathcal M_{\mathbb P}$, as $\mathbb P$ is
not a semiring. Also
$\mathrm{Id}$ uses no semiring at all, a state is just an element. The monads, on a general set $X$:

\begin{center}\small
\begin{tabular}{@{}llll@{}}
\toprule
$\mathcal M$ & semiring & $\mathcal M X$ & constraint \\
\midrule
$\mathrm{Id}$ & none & $X$ & none \\[2pt]
$\mathcal D \subset \mathcal M_{\mathbb M}$ & mass $\mathbb M$, slice $\Sigma{=}1$ & $X \to [0,1]$ & finitely many ${\ne} 0$, $\Sigma = 1$ \\[2pt]
$\mathcal T = \mathcal M_{\mathbb M}$ & mass $\mathbb M$ & $X \to \mathbb R_{\ge 0}$ & finitely many ${\ne} 0$ \\[2pt]
$\log\mathcal T = \mathcal M_{\mathbb L}$ & log $\mathbb L$ & $X \to \mathbb R \cup \{-\infty\}$ & finitely many ${\ne} {-\infty}$ \\
\bottomrule
\end{tabular}
\end{center}

\paragraph{Truth spaces.} Now apply each monad to the boolean to get the four truth spaces:

\begin{center}\small
\begin{tabular}{@{}llll@{}}
$\mathcal M$ & $\mathcal M\mathbb B$ & $\cong$ & read out \\
\midrule
$\mathrm{Id}$ & $\mathbb B$ & $=$ & $\mathbb B$ \\[2pt]
$\mathcal D$ & $\mathbb B \to [0,1],\ \Sigma {=} 1$ & $p \mapsto p(1)$ & $[0,1]$ \\[2pt]
$\mathcal T$ & $\mathbb B \to \mathbb R_{\ge 0}$ & $w \mapsto (w(0),\ w(1))$ & $\mathbb R_{\ge 0}^{2}$ \\[2pt]
$\log\mathcal T$ & $\mathbb B \to \mathbb R \cup \{-\infty\}$ & $a \mapsto (a(0),\ a(1))$ & $(\mathbb R {\cup} \{-\infty\})^{2}$ \\
\end{tabular}
\end{center}

\noindent
What the
isomorphism does differs by row. For
$\mathcal D$ it forgets the redundant coefficient,
$p(0) = 1 - p(1)$, so one number suffices.
For $\mathcal T$ and $\log\mathcal T$ nothing is redundant, so
nothing can be forgotten: the readout is just the pair of coefficients and the space stays two dimensional.

\paragraph{Lifted connectives.} For $\mathbb B$ the truth space is the
scalar carrier itself: its operations are the ones being lifted. For
$[0,1] = \mathcal D\mathbb B$ the readout is again scalar and the lift
\emph{derives} the $[0,1]$ row: binding two independent truth values and
pushing forward along the Boolean operation gives
$p \otimes q = pq$, $p \oplus q = 1-(1-p)(1-q) = p+q-pq$ and
$\neg p = 1-p$; the $\vee$, $\wedge$ columns ($\max$, $\min$) are the
order family, induced by the order, not the monad. For the squares
nothing collapses as an element is a pair:

\begin{center}\footnotesize\setlength{\tabcolsep}{3pt}
\begin{tabular}{@{}l|llll|llll|l@{}}
$\mathcal M\mathbb B$ & $\bot$ & $\top$ & $\vee$ & $\wedge$ & $0$ & $1$ & $\oplus$ & $\otimes$ & $\neg$ \\
\hline
$\mathbb B$ & $0$ & $1$ & $\mathtt{or}$ & $\mathtt{and}$ & $0$ & $1$ & $\mathtt{or}$ & $\mathtt{and}$ & $\mathtt{not}$ \\[2pt]
$[0,1]$ & $0$ & $1$ & $\max$ & $\min$ & $0$ & $1$ & $p {+} q {-} p\,q$ & $p\,q$ & $1 - p$ \\[2pt]
$\mathbb R_{\ge 0}^{2}$ & $(\infty, 0)$ & $(0, \infty)$ & join & meet & $(0,1)$ & $(1,0)$ & $\oplus$ & $\otimes$ & swap \\[2pt]
$(\mathbb R {\cup} \{-\infty\})^{2}$ & $(\infty, -\infty)$ & $(-\infty, \infty)$ & join & meet & $(0, -\infty)$ & $(-\infty, 0)$ & $\oplus$ & $\otimes$ & swap \\
\end{tabular}
\end{center}

\noindent
The lifted $\otimes$ and $\oplus$ come from one mechanism: form the outer product
matrix $(w_i\,v_j)_{i,j \in \mathbb B}$, all four pairwise products;
then regroup its entries along the Boolean operation.
For $\otimes$, the entry with $i \wedge j = 1$, that is $w_1 v_1$, is
slot $1$, and the three entries with $i \wedge j = 0$ are summed into
slot $0$; for $\oplus$ dually. For the mass semiring in
$\mathbb R_{\ge 0}^{2}$ the connectives are:
\begin{align*}
  w \otimes v &= \bigl(w_0 v_0 + w_0 v_1 + w_1 v_0,\ \ w_1 v_1\bigr),
  &
  w \oplus v &= \bigl(w_0 v_0,\ \ w_0 v_1 + w_1 v_0 + w_1 v_1\bigr),\\
  w \wedge v &= \bigl(\max(w_0, v_0),\ \min(w_1, v_1)\bigr),
  &
  w \vee v &= \bigl(\min(w_0, v_0),\ \max(w_1, v_1)\bigr),\\
  \neg w &= (w_1,\ w_0).
\end{align*}
In log
coordinates the domains and codomains are the same over
$(\mathbb R \cup \{-\infty\})^{2}$, products become sums, sums become
$\mathrm{lse}$:
\begin{align*}
  a \otimes b &= \bigl(\mathrm{lse}(a_0{+}b_0,\ a_0{+}b_1,\ a_1{+}b_0),\ \ a_1 {+} b_1\bigr),
  &
  a \oplus b &= \bigl(a_0 {+} b_0,\ \ \mathrm{lse}(a_0{+}b_1,\ a_1{+}b_0,\ a_1{+}b_1)\bigr),\\
  a \wedge b &= \bigl(\max(a_0, b_0),\ \min(a_1, b_1)\bigr),
  &
  a \vee b &= \bigl(\min(a_0, b_0),\ \max(a_1, b_1)\bigr),\\
  \neg a &= (a_1,\ a_0).
\end{align*}
The order behind $\wedge$, $\vee$ is componentwise with slot $0$
reversed: $w \le v$ iff $w_1 \le v_1$ and $w_0 \ge v_0$; $\wedge$,
$\vee$ are its meet and join, and the bounds $\bot = (\infty, 0)$,
$\top = (0, \infty)$ exist after adjoining $\infty$ to the carrier.

\paragraph{Readings.} Only now the interpretation. Slot $0$ carries the
weight against, slot $1$ the weight for. The mass is
$Z(w) = w_0 + w_1$, and it is multiplicative,
$Z(w \otimes v) = Z(w)\,Z(v)$. Under this reading, $\otimes$ says: the
weight for $w \otimes v$ is $w_1 v_1$, both for, independent weights
multiply, and the weight against sums the three ways at least one is
against. The order says: more weight for, less against. The $[0,1]$ row
is the mass one slice, $Z(w) = 1$, where one number suffices; that is
the normalization step from $\mathcal T$ to $\mathcal D$. The finer
reading of a pair, the split into log odds $\ell = a_1 - a_0$ and log
mass $\log Z$, is taken up in the coordinates paragraph below.

\paragraph{What lifts and what does not.} The precise question: take
the lattice semiring $\mathbb B$, and apply to it a monad that
itself comes from a lattice semiring $S$ (or $\mathrm{Id}$,
trivially). Is the result a lattice semiring again, and which
part of its structure comes from where? Three lemmas, then the theorem.

\begin{lemma}[pointwise]
$\mathcal M_S(\mathbb B) = S^{\mathbb B}$ carries all scalar operations
and the order componentwise, and a power of a lattice semiring
is again one.
\emph{Proof:} products in the variety, componentwise. \qed
\end{lemma}

\begin{lemma}[linear laws lift]
A commutative monad is symmetric monoidal
(\citealp[Thm.~3.2]{kockMonadsSymmetric1970};
\citealp[Thm.~2.3]{kockStrongFunctors1972}), and $\mathcal M_S$ is
commutative iff $S$ is \citep[Lem.~23]{coumansJacobsScalars2013}. Hence
every equation in which each variable occurs exactly once on each side
is preserved by the lift: associativity, commutativity, units, the
involution, De Morgan. In particular
$(\mathcal M\mathbb B, \oplus, \eta(0))$ and
$(\mathcal M\mathbb B, \otimes, \eta(1))$ are commutative monoids and
$\neg$ is a De Morgan involution between them.
\end{lemma}

\begin{lemma}[copying laws do not lift]
Equations that duplicate a variable fail in general: idempotence,
absorption, lattice distributivity.
\emph{Proof:} in $\mathcal D\mathbb B$,
$p \otimes p = p^{2} \ne p$; a lifted operation reads its arguments as
independent. \qed\\
This failure is wanted: non idempotence is what makes truth graded.
\end{lemma}

\begin{theorem}[three layers of lattice semirings]
Let $S$ be a lattice semiring, and apply
$\mathcal M_S$ to the lattice semiring
$\mathbb B$. Then $\mathcal M_S\mathbb B$ is a lattice semiring
again:
\begin{itemize}
\item[(i)] \emph{pointwise}, $S^{\mathbb B}$ is a lattice semiring, so the construction closes and iterates:
$\mathcal M_S(\mathbb B)$ can serve as the weight semiring of a next
layer (semiring rows only: $S \in \{\mathbb B, \mathbb M, \mathbb L\}$,
not $\mathbb P$);
\item[(ii)] the \emph{lifted} family $\otimes$, $\oplus$, $\neg$ with
units $\eta(0)$, $\eta(1)$ satisfies every linear Boolean law and no
copying law: two commutative monoids with a De Morgan involution, not a
semiring;
\item[(iii)] the \emph{order} family $\wedge$, $\vee$, $\bot$, $\top$
is induced by the order of $S$, componentwise with the $0$ slot
reversed, not by the monad; the bounds exist after adjoining $\infty$.
\end{itemize}
The copying axioms are exactly the order generating ones ($u \le v$ iff
$u \wedge v = u$), so (ii) and (iii) are disjoint contributions: both
families are logic, differing only in origin. They coincide exactly
where the carrier is idempotent, on $\mathbb B$ itself: G\"odel and
product connectives side by side, collapsing only there. The units
$\eta(0)$, $\eta(1)$ agree with the bounds $\bot$, $\top$ only in the
normalized rows.
\emph{Proof:} (i) is the first lemma, (ii) the second and third lemma,
(iii) the construction of the order on pairs above. \qed
\end{theorem}

\paragraph{Four monads, one basis.} All four monads on the
bits ($\mathbb B = \{0, 1\}$) share the same basis: the computational basis of information theory\footnote{Also called the Z basis in
quantum information theory \citep{heunenCategoriesQuantumTheory2019}.}, the two kets $\ket{0}$ and $\ket{1}$,
false as $0$ and true as $1$. What varies are the coefficients and the constraint:

\begin{center}\small
\begin{tabular}{@{}llll@{}}
monad & a state on $\mathbb B$ is & coefficients & constraint \\
\midrule
$\mathrm{Id}$ & $b_0\,\ket{0} + b_1\,\ket{1}$ & $b_0, b_1 \in \{0,1\}$ & $b_0 + b_1 = 1$ \\[2pt]
$\mathcal D$ & $p_0\,\ket{0} + p_1\,\ket{1}$ & $p_0, p_1 \in [0,1]$ & $p_0 + p_1 = 1$ \\[2pt]
$\mathcal T$ & $w_0\,\ket{0} + w_1\,\ket{1}$ & $w_0, w_1 \in \mathbb R_{\ge 0}$ & none \\[2pt]
$\log\mathcal T$ & $a_0\,\ket{0} + a_1\,\ket{1}$ & $a_0, a_1 \in \log \mathbb R_{\ge 0}$ & none
\end{tabular}
\end{center}
\noindent
A classical truth value \emph{is} one of the two bits. In
$\mathcal D$ it is a probability distribution over the two bits. In
$\mathcal T$ the normalization is dropped, allowing any two nonnegative masses. In
$\log\mathcal T$ the same states are written in log coordinates,
$a_0 = \log w_0$ and $a_1 = \log w_1$. The bits: in $\mathcal T$ the bit
$\ket{0}$ is the basis vector $(1, 0)$; taking logs componentwise
($\log 1 = 0$, $\log 0 = -\infty$) it becomes $(0, -\infty)$. On the other hand, the bit $\ket{1}$ is the vector $(0, 1)$; taking logs $(-\infty, 0)$. As in Dirac
notation, the ket symbol is representation free: $\ket{0}$ names the
same basis state in every row, and only its coordinate vector changes
with the carrier, since the units correspond under the isomorphism,
$\eta_{\log} = \log \circ \eta$. When the representation matters we
write the coordinates, as in the table; never $\ket{\log 0}$, because
$\log 0$ is not a label of $\mathbb B$.


\paragraph{Coefficients, mass, matrices.} From here on we write
$\braket{u}{w} := w(u)$ for the coefficient of $\ket{w}$ at $u$. The costate
$\bra{\Sigma} := \bra{0} + \bra{1}$ sums them: the mass is
$Z(w) = \braket{\Sigma}{w}$, and in log coordinates
$\log Z(a) = \mathrm{lse}\bigl(\braket{0}{a}, \braket{1}{a}\bigr)$.
A lifted pure map is the matrix $\sum_x \ketbra{h(x)}{x}$ with
exactly one $1$ per column, so it preserves mass: every column sums to
one. Negation is the $X$ gate,
$\neg = \ketbra{0}{1} + \ketbra{1}{0}$.

\paragraph{The chain of maps.} One step at a time, each map typed:
\[
  \mathbb B
  \;\xrightarrow{\ \eta_{\mathcal D}\ }\;
  \mathcal D\mathbb B
  \;\xrightarrow{\ \iota\ }\;
  \mathcal T\mathbb B
  \;\xrightarrow[\cong]{\ \log\ }\;
  \log\mathcal T\mathbb B
\]
\begin{itemize}
\item $\eta_{\mathcal D} \colon \mathbb B \to \mathcal D\mathbb B$,
  $u \mapsto \ket{u}$: the point mass is the ket, Dirac in both senses
  of the word. Classical truth embeds as the two extreme points of the
  interval. Injective, far from surjective: $\mathcal D$ adds the whole
  interval of graded truth. No canonical map back; decisions are argmax.
\item $\iota \colon \mathcal D\mathbb B \to \mathcal T\mathbb B$: the
  inclusion of the mass one slice $\braket{\Sigma}{w} = 1$. A
  section of normalization
  $\ket{w} \mapsto \ket{w} / \braket{\Sigma}{w}$, not an
  isomorphism: $\mathcal T$ adds exactly one dimension, the mass.
\item $\log \colon \mathcal T\mathbb B \to \log\mathcal T\mathbb B$,
  coordinatewise, inverse $\exp$: a genuine isomorphism, the semiring
  isomorphism $(\mathbb R_{\ge 0}, +, \cdot) \cong
  (\mathbb R \cup \{-\infty\}, \mathrm{lse}, +)$ applied to the carrier.
  So $\mathcal T\mathbb B$ and $\log\mathcal T\mathbb B$ are the same
  space in different coordinates, not two different spaces; we always
  compute in the log coordinates, where products become sums,
  marginalisation becomes $\mathrm{lse}$, and the numerics are stable. In
  log coordinates the ket $\ket{u}$ becomes the one hot vector with
  $0$ at $u$ and $-\infty$ elsewhere.
\end{itemize}
In this notation the bridge maps factor as
$\mathrm{enc} = \iota \seq \log$ and
$\mathrm{dec} = \exp \seq \bigl(\ket{w} \mapsto \ket{w} / \braket{\Sigma}{w}\bigr)$, and the mass one slice in log coordinates is
$\{\mathrm{lse}(a) = 0\}$, the image of $\mathrm{nrm}$.

\paragraph{The Boolean case, in coordinates.} The shift decomposition
splits $\log\mathcal T\mathbb B$ into two independent coordinates,
\[
  \ket{a} \;\cong\; (\ell,\ \log Z),
  \qquad
  \ell := \braket{1}{a} - \braket{0}{a},
  \quad
  \log Z := \mathrm{lse}\bigl(\braket{0}{a},\ \braket{1}{a}\bigr),
\]
an isomorphism, since $\braket{1}{a} = \log Z + \log\sigma(\ell)$
and $\braket{0}{a} = \log Z + \log\sigma(-\ell)$ with $\sigma$ the
logistic function. In these coordinates:
\begin{itemize}
\item $\mathrm{dec}$ is the sigmoid: softmax on two elements is
  $\sigma(\braket{1}{a} - \braket{0}{a})$, so $\mathrm{dec}$
  projects to $\ell$ and applies $\sigma$. The probability of $1$ lives
  entirely in $\ell$.
\item The shift acts only on $\log Z$: adding $c\,\mathbf 1$ leaves
  $\ell$ unchanged. So $\ell$ is the semantics and $\log Z$ is the gauge;
  a Boolean truth value factors as (meaning, evidence mass).
\item $\mathrm{nrm}$ sets $\log Z := 0$ and keeps $\ell$. The mass one
  slice consists of the log probability states with
  $\braket{0}{a} = \log(1{-}p)$, $\braket{1}{a} = \log p$, and
  the coefficient
  $\braket{1}{\mathrm{nrm}(a)} = \log p = -\,\mathrm{softplus}(-\ell)
  \in [-\infty, 0]$ is exactly what the knowledge loss
  $-\log \varphi_\theta$ consumes.
\item A standard single logit binary head works in the quotient by the
  shift: it parametrizes $\ell$ and never materializes $\log Z$. The
  classical logit of logistic regression is the shift invariant part of
  the Boolean log truth value.
\end{itemize}
Compare the $\mathcal D$ side: $\mathcal D(\mathbb B) \cong [0,1]$, a
single probability (the remark at l.155 to 156). In $\mathcal T$ that
collapse is not available: the truth value is genuinely two dimensional,
$[0,1]$ worth of meaning plus one real of mass, and the extra dimension
is exactly the tilt coordinate of this analysis. The implementation still
keeps both coefficients rather than $\ell$ alone because the lifted
connectives act linearly on states, as the ket bra matrices above, while
on $\ell$ alone every connective would need its own $\mathrm{lse}$
correction. Two coefficients buy linear algebraic semantics; the price is
carrying the gauge, which is the dec commutation question in miniature.

\paragraph{Conjunction, typed.} For $\Omega = \mathbb B$ write
$\mathcal T\Omega = (\mathbb R \cup \{-\infty\})^{2}$ in log coordinates
and $\mathcal D\Omega = \{\braket{\Sigma}{p} = 1\}$. The lifted
conjunction factors through the pairing of the two effectful wires
($\eta$ followed by the two strength binds of Definition~10), written
$\mathrm{dst}$; the pair basis kets are $\ket{uv}$, and in log space
the outer product is the outer sum. All arrows with their domains and
codomains:
\begin{align*}
  \wedge \;&\colon\; \mathbb B \otimes \mathbb B \to \mathbb B,
    && u \wedge v = uv,\ \text{only } \ket{11} \mapsto \ket{1},\\
  \mathrm{dst}_{\mathcal T} \;&\colon\; \mathcal T\Omega \otimes \mathcal T\Omega \to \mathcal T(\Omega \otimes \Omega),
    && \braket{uv}{a \oplus b} = \braket{u}{a} + \braket{v}{b},\\
  \mathcal T(\wedge) \;&\colon\; \mathcal T(\Omega \otimes \Omega) \to \mathcal T\Omega,
    && \ketbra{1}{11} + \ket{0}\bigl(\bra{00} + \bra{01} + \bra{10}\bigr),\\
  \mathrm{dec}_\Omega \;&\colon\; \mathcal T\Omega \to \mathcal D\Omega,
    && \braket{u}{\mathrm{dec}(a)} = e^{\braket{u}{a}} / Z(a),\ \
       \braket{1}{\mathrm{dec}(a)} = \sigma(\ell).
\end{align*}
The conjunctions of the two readings are the composites
\begin{align*}
  \otimes_{\mathcal T} \;&:=\; \mathrm{dst}_{\mathcal T} \seq \mathcal T(\wedge)
    \;\colon\; \mathcal T\Omega \otimes \mathcal T\Omega \to \mathcal T\Omega,\\
  \otimes_{\mathcal D} \;&:=\; \mathrm{dst}_{\mathcal D} \seq \mathcal D(\wedge)
    \;\colon\; \mathcal D\Omega \otimes \mathcal D\Omega \to \mathcal D\Omega,
\end{align*}
explicitly, in full,
\[
  \braket{1\,}{\,a \otimes_{\mathcal T} b}
  = \braket{1}{a} + \braket{1}{b},
  \qquad
  \braket{0\,}{\,a \otimes_{\mathcal T} b}
  = \mathrm{lse}\bigl(\braket{0}{a}{+}\braket{0}{b},\
    \braket{0}{a}{+}\braket{1}{b},\
    \braket{1}{a}{+}\braket{0}{b}\bigr),
\]
\[
  \braket{1\,}{\,p \otimes_{\mathcal D} q}
  = \braket{1}{p}\,\braket{1}{q},
  \qquad
  \braket{0\,}{\,p \otimes_{\mathcal D} q}
  = 1 - \braket{1}{p}\,\braket{1}{q},
\]
so $\otimes_{\mathcal D}$ is the product t norm. How
$\otimes_{\mathcal T}$ composes with softmax down to $\mathcal D\Omega$:
\[\begin{tikzcd}[ampersand replacement=\&, column sep=3.0em]
  {\mathcal T\Omega \otimes \mathcal T\Omega} \& {\mathcal T(\Omega \otimes \Omega)} \& {\mathcal T\Omega} \\
  {\mathcal D\Omega \otimes \mathcal D\Omega} \& {\mathcal D(\Omega \otimes \Omega)} \& {\mathcal D\Omega}
  \arrow["{\mathrm{dst}_{\mathcal T}}", from=1-1, to=1-2]
  \arrow["{\mathcal T(\wedge)}", from=1-2, to=1-3]
  \arrow["{\mathrm{dec} \otimes \mathrm{dec}}"', from=1-1, to=2-1]
  \arrow["{\mathrm{dec}_{\Omega \otimes \Omega}}"{description}, from=1-2, to=2-2]
  \arrow["{\mathrm{dec}_\Omega}", from=1-3, to=2-3]
  \arrow["{\mathrm{dst}_{\mathcal D}}"', from=2-1, to=2-2]
  \arrow["{\mathcal D(\wedge)}"', from=2-2, to=2-3]
\end{tikzcd}\]
The left square commutes by the tensor lemma ($\mathrm{dec}$ is
monoidal), the right square by the pure map lemma ($\mathrm{dec}$ is
natural in ket bra matrices), hence the outer rectangle:
$\mathrm{dec}(a \otimes_{\mathcal T} b) =
 \mathrm{dec}(a) \otimes_{\mathcal D} \mathrm{dec}(b)$, always. Chasing
the $\ket{1}$ coefficient along both paths shows why:
\begin{align*}
  (\ket{a}, \ket{b})
  \;&\stackrel{\mathrm{dst}_{\mathcal T}}{\longmapsto}\; \ket{a \oplus b}
  \;\stackrel{\mathcal T(\wedge)}{\longmapsto}\;
    \braket{1}{{\cdot}} = \braket{1}{a} + \braket{1}{b}
  \;\stackrel{\mathrm{dec}}{\longmapsto}\;
    \frac{e^{\braket{1}{a} + \braket{1}{b}}}{Z(a)\,Z(b)},\\
  (\ket{a}, \ket{b})
  \;&\stackrel{\mathrm{dec} \otimes \mathrm{dec}}{\longmapsto}\;
    \bigl(\mathrm{dec}(a), \mathrm{dec}(b)\bigr)
  \;\stackrel{\otimes_{\mathcal D}}{\longmapsto}\;
    \frac{e^{\braket{1}{a}}}{Z(a)} \cdot
    \frac{e^{\braket{1}{b}}}{Z(b)},
\end{align*}
equal because the pushforward preserves mass:
$Z\bigl(\mathcal T(\wedge)(a \oplus b)\bigr) = Z(a \oplus b) =
Z(a)\,Z(b)$. The $\ket{0}$ coefficient is the complement in both
readings.

In the three coordinate systems above:
\begin{itemize}
\item In ($\ket{1}$ coefficient, mass) coordinates
  $\otimes_{\mathcal T}$ is componentwise addition of logs: the
  $\braket{1}{{\cdot}}$ log weights add and the log masses add; the
  $\ket{0}$ coefficient is bookkeeping,
  $\braket{0}{w} = Z - \braket{1}{w}$.
\item On the mass one slice it is the product t norm:
  $p \otimes q = pq$, dually $p \oplus q = p + q - pq$ (the $\ket{0}$
  coefficients multiply), and negation is the $X$ gate, so De Morgan
  holds by construction. The fuzzy product semantics is the shadow of
  probabilistic conjunction of independent events (usable in the
  Appendix B narrative of T2).
\item On the log odds alone,
  $\ell_{a \otimes b} = -\,\mathrm{lse}(-\ell_a,\ -\ell_b,\
  -\ell_a - \ell_b)$; sanity check: two fair coins, $\ell_a = \ell_b = 0$,
  give $-\log 3$, odds $1{:}3$ for probability $\tfrac14$ of $\ket{1}$.
  On $\ell$ every connective needs its own nonlinear formula, while on
  full states everything stays outer product then pushforward.
\end{itemize}

\section{Categorical Batching}

\begin{proposition}[Pointwise evaluation]
  \label{prop:pointwise}
  For each $i \in \underline B$, evaluation $\mathrm{ev}_i \colon \mathcal B\mathcal M\,X \to
  \mathcal M\,X$, $m \mapsto m(i)$, is a monad morphism, and therefore commutes with the
  interpretation of \doNotation programs: for a batch $s \colon \underline B \to S$,
  \[
    \sem{\varphi}^{\mathcal B\mathcal M}(s)(i) \;=\; \sem{\varphi}^{\mathcal M}(s_i)
    \qquad (i \in \underline B).
  \]
  \end{proposition}
  
  \noindent
  One batched run thus yields all per-sample truth values; over the full
  sample these are the $N$ numbers $\widehat L$ averages, a mini-batch gives an unbiased
  estimate of $\widehat L$.

% ============================================================
% REVISION TODOS (NeSy26 reviews, 2026-08-05). Working notes.
% REMOVE THIS ENTIRE SECTION BEFORE RESUBMISSION.
% Tn labels match the task titles in the Admination "Nesy2026" process.
% Source line references (l.NNN) refer to commit 58b317f and will drift
% as the file is edited; the \ref anchors are stable.
% ============================================================
\newpage
\begingroup
\emergencystretch=2.5em
\providecommand{\qed}{\hfill$\square$}
\section*{Revision TODOs (NeSy26 reviews): remove before resubmission}

\subsection*{Daniel}

\begin{description}

\item[T1: answer the only MNIST point (Rev.~2).]
Three steps. (i) For the rebuttal: $N{=}4$ already is the standard hard
case. It needs $8$ CNN calls and a joint space of $10^{8}$ digit
combinations per sample, all summed out exactly. DeepProbLog and LTN time
out there (Table~\ref{tab:addition}), and A-NeSI \citep{kriekenANeSI2023}
exists exactly because this size was considered out of reach for exact
inference. (ii) Tone down two spots that promise more than the experiments
show: the abstract (``we stay in a uniform framework that applies to many
first-order NeSy approaches'', l.71) and the Conclusion (l.620). Say the
framework \emph{supports} the other semantics instead of suggesting we
tested them all. (iii) Once T2 is in, the answer to ``only one benchmark''
is: we claim breadth across semantics, not across datasets.

\item[T2: fuzzy experiment, {\L}ukasiewicz logic (new Appendix B).]
Run one fuzzy example: identity monad, truth space $\Omega=[0,1]$,
{\L}ukasiewicz connectives, $p$-mean quantifiers (the LTN setup). Take the
smokers friends cancer example, \S4.7 of
\citet{DBLP:journals/ai/BadreddineGSS22} (checked: it is there; knowledge
base completion, $14$ constants, predicates $S,F,C$). It is a good pick for
two reasons: the knowledge base cannot be fully satisfied, so fuzzy truth
values are the natural reading, and it has no grammar or derivation
structure, so DeepStochLog cannot express it (used again in T11). Code: one
new \texttt{Fuzzy} monad instance next to \texttt{LogTens} and
\texttt{Dist}; the monad polymorphic formula code of Sec.~\ref{sec:mnist}
stays as is. Paper: new Appendix B, about one page (axioms, one results
table), referenced from Sec.~\ref{sec:eval} and Sec.~\ref{sec:related}.

\item[T3: can $\mathcal D$ and $\mathcal T$ disagree on a formula? (Rev.~3).]
They agree for programs shaped
like our MNIST axiom, and they disagree in general. The shape of the monad both programs run in is:
\[
  \mathcal M(X) \;=\; \{\, f : X \to S \,\},
\]
finitely supported $S$ valued maps; for $\mathcal T$ the semiring $S$ is
the log semiring, for $\mathcal D$ the semiring $S$ is the probability semiring. Next the bridge definition:
\[
  \mathrm{digit}^{\mathcal D}_\theta \;:=\;
  \mathrm{enc} \seq \mathrm{digit}^{\mathcal T}_\theta \seq  \mathrm{dec}.
\]
The question is wether the following equation holds:
\begin{align*}
  &\mathrm{dec} \,\circ\, \DO_{\mathcal T}\{\,
    d_1 \leftarrow (\mathrm{digit}^{\mathcal T}_\theta \circ \mathrm{enc})(x);\
    d_2 \leftarrow (\mathrm{digit}^{\mathcal T}_\theta \circ \mathrm{enc})(y);\
    \mathbf{return}\,(s {=} d_1 {+} d_2) \,\}
  \\[2pt]
  &\quad\overset{?}{=}\quad
  \DO_{\mathcal D}\{\,
    d_1 \leftarrow \mathrm{digit}^{\mathcal D}_\theta(x);\
    d_2 \leftarrow \mathrm{digit}^{\mathcal D}_\theta(y);\
    \mathbf{return}\,(s {=} d_1 {+} d_2) \,\}.
\end{align*}

Write
$\seqop_{\mathcal T}$ and $\seqop_{\mathcal D}$ for the feed forward
composite,
taken with the bind of $\mathcal T$ and of $\mathcal D$ respectively, and
$\mathcal I({=}\,s)$ for the pure comparison with the observed sum. The
chains start from input distributions in
$\mathcal D(\mathsf{Image})$; the leaf is
$\mathrm{enc} \fatsemi \mathrm{dig}^{\mathcal T}_\theta \colon
\mathcal D(\mathsf{Image}) \to \mathcal T(\mathsf{Digit})$. The left side
of the equation is the chain
\[
  \mathrm{LHS} \;=\;
  \bigl(
    (\mathrm{enc} \fatsemi \mathrm{dig}^{\mathcal T}_\theta)
    \otimes
    (\mathrm{enc} \fatsemi \mathrm{dig}^{\mathcal T}_\theta)
  \bigr)
  \seqop_{\mathcal T} \mathcal I(+)
  \seqop_{\mathcal T} \mathcal I({=}\,s)
  \fatsemi \mathrm{dec},
\]
one single $\mathrm{dec}$ at the very end, and the right side is
\[
  \mathrm{RHS} \;=\;
  \bigl(
    (\mathrm{enc} \fatsemi \mathrm{dig}^{\mathcal T}_\theta \fatsemi \mathrm{dec})
    \otimes
    (\mathrm{enc} \fatsemi \mathrm{dig}^{\mathcal T}_\theta \fatsemi \mathrm{dec})
  \bigr)
  \seqop_{\mathcal D} \mathcal I(+)
  \seqop_{\mathcal D} \mathcal I({=}\,s),
\]

Every feed
forward composite unfolds into the point free chain
\[
  \mathcal I(f) \fatsemi \vec\alpha \fatsemi \eta
  \fatsemi ({\bind}\,\sigma^{(j_1)}) \fatsemi \cdots
  \fatsemi ({\bind}\,\sigma^{(j_r)})
  \fatsemi ({\bind}\,\mathcal I(g) \fatsemi \eta),
\]
So a formula is one long chain $L \fatsemi B_1 \fatsemi \cdots \fatsemi
B_r$ with $L$ the tensor of the leaves, and every factor $B_i$ is one of
four kinds: a lifted pure map (adapters, $\mathcal T(h)$), a unit $\eta$, a bind ${\bind}\,\sigma^{(j_i)}$, or a bind ${\bind}\,\mathcal I(g) \fatsemi \eta$. The question
``can $\mathrm{dec}$ be pulled out'' is therefore answered factor by
factor. We do that now, with one lemma per kind.

\paragraph{Notation:} For a finite set $X$, a $\mathcal T$ value is a log weight vector
$a \in \mathbb R^X$ (allowing $-\infty$). Define three arrow
families:
\[
  \mathrm{dec}_X(a)_x = \frac{e^{a_x}}{Z(a)}, \qquad
  \mathrm{enc}_X(p)_x = \log p(x), \qquad
  Z_X(a) = \sum_{x \in X} e^{a_x},
\]
so $\mathrm{dec}$ is softmax, $\mathrm{enc}$ is pointwise log, and $Z$ is
the total \emph{mass}, the number $\mathrm{dec}$ divides by. For
$a \in \mathcal T X$ and $k \colon X \to \mathcal T Y$ the bind of
$\mathcal T$ is $(a \bind k)_y = \mathrm{logsumexp}_x\,(a_x + k(x)_y)$; for
$p \in \mathcal D X$ and $q \colon X \to \mathcal D Y$ the bind of
$\mathcal D$ is $(p \bind q)(y) = \sum_x p(x)\, q(x)(y)$.

\begin{lemma}[pure maps]
For pure $h \colon X \to Y$:
$Z(\mathcal T(h)(a)) = Z(a)$ and
$\mathcal T(h) \fatsemi \mathrm{dec}_Y = \mathrm{dec}_X \fatsemi \mathcal D(h)$.

\emph{Proof:} The fibers $h^{-1}(y)$ split $X$ into disjoint parts, so
summing $e^{a_x}$ fiber by fiber rearranges the same total. \qed
\end{lemma}

\begin{lemma}[tensor]
$Z(a \otimes b) = Z(a)\,Z(b)$ and
$\mathrm{dec}(a \otimes b) = \mathrm{dec}(a) \otimes \mathrm{dec}(b)$.

\emph{Proof:} $e^{a_i + b_k} = e^{a_i} e^{b_k}$, then sum. \qed
\end{lemma}

Hence $\mathrm{dec}$ commutes with all the wiring of the CD
category, pairing leaves, copying observed data, permuting wires.

\begin{lemma}[units]
$\eta_{\mathcal D} \fatsemi \mathrm{enc} = \eta_{\mathcal T}$,
$\mathrm{enc} \fatsemi \mathrm{dec} = \mathrm{id}_{\mathcal D}$, hence
$\eta_{\mathcal T} \fatsemi \mathrm{dec} = \eta_{\mathcal D}$, and
$Z(\eta_{\mathcal T}(x)) = 1$.
\end{lemma}

Therefore $\mathrm{enc}$ is a section of
$\mathrm{dec}$: encoding loses nothing. The reverse composite
$\mathrm{dec} \fatsemi \mathrm{enc} \ne \mathrm{id}_{\mathcal T}$ forgets
the shift (the shift invariance remark, l.443 to 445). With the current
$\epsilon$ floor all of this holds only up to $\epsilon$ terms.

\begin{lemma}[strengths]
For the strength at $j$:
\begin{gather*}
  Z\bigl(\sigma^{(j)}(x_1,\dots,a,\dots,x_m)\bigr) \;=\; Z(a),\\
  \sigma^{(j)}_{\mathcal T} \fatsemi \mathrm{dec}
  \;=\;
  (\mathrm{id} \otimes \cdots \otimes \mathrm{dec}_j \otimes \cdots \otimes \mathrm{id})
  \fatsemi \sigma^{(j)}_{\mathcal D}.
\end{gather*}
\emph{Proof.} The plain coordinates enter as certain values of mass $1$;
apply the previous two lemmas. \qed
\end{lemma}

\noindent
So $\mathrm{dec}$ commutes with every factor of the chain except possibly
the binds. What one bind costs is exact:

\begin{lemma}[tilt]
Let $a \in \mathcal T X$ and $k \colon X \to \mathcal T Y$. Then
\[
  \mathrm{dec}(a \bind k)
  \;=\;
  \mathrm{tilt}_{k \seq Z}\bigl(\mathrm{dec}(a)\bigr) \bind (k \fatsemi \mathrm{dec}),
\]
where for a weight function $w \colon X \to \mathbb R_{>0}$
\[
  \mathrm{tilt}_w(p)(x) \;:=\; \frac{w(x)\,p(x)}{\sum_{x'} w(x')\,p(x')}.
\]
\emph{Proof.} Expand both sides; the numerators agree termwise, and the
denominator $\sum_{x,y} e^{a_x + k(x)_y} = \sum_x e^{a_x} Z(k(x))$ is
exactly the tilt's normalizer. \qed\\
\end{lemma}

In words that means that pulling $\mathrm{dec}$ through one bind costs exactly one
reweighting of the earlier distribution, by the mass of the continuation.

\begin{definition}[mass preserving]
An arrow $k \colon X \to \mathcal T Y$ is \emph{mass preserving} if
$x \mapsto Z(k(x))$ is constant. 
\end{definition}

By the lemmas above this class contains
$h \fatsemi \eta_{\mathcal T}$ for pure $h$, all lifts $\mathcal T(h)$, the
units, the strengths, and it is closed under $\fatsemi$, $\otimes$ and the
wiring. A neural leaf is generally \emph{not} in the class: its logits
carry whatever mass they like.

\begin{theorem}[pull out]
Let $\Phi = L \fatsemi B_1 \fatsemi \cdots \fatsemi B_r$ be the point free
chain of a feed forward composite, $L$ the tensor of the leaves, each $B_i$
a lifted pure map, a unit, a strength factor, or a bind ${\bind}\,k_i$. If
every bound continuation $k_i$ is mass preserving, then
\[
  \Phi \fatsemi \mathrm{dec}
  \;=\;
  \bigl(L \fatsemi (\mathrm{dec} \otimes \cdots \otimes \mathrm{dec})\bigr)
  \fatsemi B_1^{\mathcal D} \fatsemi \cdots \fatsemi B_r^{\mathcal D},
\]
the same chain read in $\mathcal D$, with one $\mathrm{dec}$ per leaf. If
some $k_i$ has non constant mass, the two sides differ in general, and the
tilt formula gives the exact discrepancy: one $\mathrm{tilt}_{Z \circ k_i}$
per offending bind.

\emph{Proof.} Induction over $r$: each factor commutes with $\mathrm{dec}$
by the four factor lemmas; a bind factor commutes by the tilt formula,
because $\mathrm{tilt}_w = \mathrm{id}$ exactly when $w$ is constant. The
converse is the example below. \qed
\end{theorem}

\begin{corollary}[how wrong can it get]
If $m_i \le Z(k_i(x)) \le M_i$ for all $x$, then the two truth values
satisfy
$\bigl|\log t_{\mathcal T} - \log t_{\mathcal D}\bigr|
 \le \sum_i \log (M_i / m_i)$,
and the relation is the identity exactly when every $M_i = m_i$. Proved
above for one bind (both sides average the same values under densities
that differ by at most $M/m$); the chain version is the remaining piece of
Till's monotonicity requirement, still to do.
\end{corollary}

\begin{example}[failure, exact numbers]
$X = Y = \{0,1\}$, $a = (0,0)$, $k(0) = (0,0)$, $k(1) = (\ln 100,\, 0)$.
Masses: $Z(k(0)) = 2$, $Z(k(1)) = 101$. The tilt turns
$\mathrm{dec}(a) = (\tfrac12, \tfrac12)$ into
$(\tfrac{2}{103}, \tfrac{101}{103})$, so
$\mathrm{dec}(a \bind k)(0) = 101/103 \approx 0.981$, while the untilted
$\mathcal D$ side gives
$\tfrac14 + \tfrac{50}{101} \approx 0.745$. One bind, one tilt, visibly
different.
\end{example}

\begin{example}[the MNIST chains commute]
In LHS and RHS above, the leaves are
$\mathrm{enc} \fatsemi \mathrm{dig}^{\mathcal T}_\theta$ on arbitrary input
distributions, and after them every bind binds a strength, or
$\mathcal I(+) \fatsemi \eta$, or $\mathcal I({=}\,s) \fatsemi \eta$: all
mass preserving. The theorem applies, so $\mathrm{LHS} = \mathrm{RHS}$: the
board equation is true for the MNIST axiom. Concretely the theorem
specializes to the convolution identity
$\mathrm{dec}(\mathrm{logconv}(a,b)) = \mathrm{dec}(a) * \mathrm{dec}(b)$,
which we also checked over $20{,}000$ random trials. In the four wire
cases of
Definition~9 of the theory paper: case one (plain into plain) and case two
(plain into circled, adapter $\eta$) always commute with $\mathrm{dec}$
(the pure map and unit lemmas); case three (circled into circled,
computation passed as data) merges nothing, it hands the question to the
consuming symbol, which is exactly where the two sided digit and the
bridge live; case four (circled into plain) is a bind and commutes exactly
when its continuation is mass preserving (the tilt lemma). So the safe
fragment reads: every case four wire after the neural leaves has a mass
preserving continuation, which holds when only plain output symbols occur
downstream of the leaves. It fails as soon as a computation whose mass
depends on an earlier bound value is bound again by a later case four
wire.
\end{example}

\paragraph{What standard machine learning does, in our arrows.} Ordinary
deep learning also trains on logits, but it normalizes every head before
anything downstream can consume it. That discipline is an arrow we already
know:

\begin{definition}[normalizer]
$\mathrm{nrm}_X := \mathrm{dec}_X \seq \mathrm{enc}_X \colon
\mathcal T X \to \mathcal T X$, concretely
$\mathrm{nrm}(a) = a - \log Z(a)\,\mathbf 1$, the log softmax. It
satisfies (i) $\mathrm{nrm} \seq \mathrm{dec} = \mathrm{dec}$
(predictions unchanged), (ii) $Z(\mathrm{nrm}(a)) = 1$ (mass one),
(iii) $\mathrm{nrm}(a + c\,\mathbf 1) = \mathrm{nrm}(a)$ (kills the
shift), (iv) $\mathrm{nrm} \seq \mathrm{nrm} = \mathrm{nrm}$.
\end{definition}

So the composite $\mathrm{dec} \seq \mathrm{enc}$ of the units lemma, the
one that fails to be the identity, is not a defect: it \emph{is} the
normalized head. What it forgets (mass, shift) is exactly what standard
machine learning throws away after every block of logits. As everywhere,
$-\infty$ entries need the log semiring of T6.
\[\begin{tikzcd}[ampersand replacement=\&]
  {\mathcal T X} \&\& {\mathcal T X} \\
  \& {\mathcal D X}
  \arrow["{\mathrm{nrm} \,=\, \mathrm{dec} \seq \mathrm{enc}}", from=1-1, to=1-3]
  \arrow["{\mathrm{dec}}"', from=1-1, to=2-2]
  \arrow["{\mathrm{enc}}"{description}, from=2-2, to=1-3]
  \arrow["{\mathrm{dec}}", from=1-3, to=2-2]
\end{tikzcd}\]
Standard supervised classification is the one leaf chain
\[
  \mathrm{net} \seq \mathrm{nrm}
  \qquad\text{with loss}\qquad
  \ell(a, y) \;=\; -\,\mathrm{nrm}(a)_y \;=\; -\log \mathrm{dec}(a)_y,
\]
manifestly a function of the decoded distribution only; by (iii) the
shift direction carries zero gradient. That is why training on logits is
harmless there: the objective factors through $\mathrm{dec}$ by
construction. In our language, after a $\mathrm{nrm}$ everything
downstream sees only mass one inputs, every continuation is mass
preserving, and the pull out theorem holds trivially. Standard machine
learning never meets the tilt because it normalizes before anything can
consume raw mass.

\paragraph{What we do, in the same picture.} \nesycattorch\ removes the
$\mathrm{nrm}$'s and evaluates the whole program on raw logits, decoding
once at the boundary (``score now, normalise once at the boundary'',
l.443 to 445). For the MNIST axiom the ladder is
\[\begin{tikzcd}[ampersand replacement=\&, column sep=2.4em]
  {\mathcal T(\mathsf{Digit}) \otimes \mathcal T(\mathsf{Digit})} \&\&
  {\mathcal T(\mathsf{Nat})} \&\& {\mathcal T\,\Omega} \\
  {\mathcal D(\mathsf{Digit}) \otimes \mathcal D(\mathsf{Digit})} \&\&
  {\mathcal D(\mathsf{Nat})} \&\& {\mathcal D\,\Omega}
  \arrow["{\text{binds},\ \mathcal I(+)}", from=1-1, to=1-3]
  \arrow["{\mathcal I({=}\,s)}", from=1-3, to=1-5]
  \arrow["{\mathrm{dec} \otimes \mathrm{dec}}"', from=1-1, to=2-1]
  \arrow["{\mathrm{dec}}"{description}, from=1-3, to=2-3]
  \arrow["{\mathrm{dec}}", from=1-5, to=2-5]
  \arrow["{\text{binds},\ \mathcal I(+)}"', from=2-1, to=2-3]
  \arrow["{\mathcal I({=}\,s)}"', from=2-3, to=2-5]
\end{tikzcd}\]
and the pull out theorem is the statement that $\mathrm{dec}$ can walk
from the right edge to the left edge, one square at a time. Here both
squares commute (pure continuations: the pure map, tensor and strength
lemmas), so deferring \emph{all} normalization is free and the
implemented loss is exactly the likelihood. In a chained program one top
edge is a bind whose continuation is itself neural; that square fails by
exactly $\mathrm{tilt}_{k \seq Z}$ and $\mathrm{dec}$ gets stuck there.

\paragraph{The exact parallel of the machine learning discipline.} To
restore the walk for arbitrary programs, do per head what standard
machine learning does everywhere: post compose every neural symbol with
the normalizer,
$\widehat{\mathrm{dig}}{}^{\mathcal T}_\theta
 := \mathrm{dig}^{\mathcal T}_\theta \seq \mathrm{nrm}$.

\begin{corollary}[normalized heads]
If every neural symbol interpretation is post composed with
$\mathrm{nrm}$, every leaf and every bound continuation has mass one,
hence is mass preserving, and the pull out theorem applies to
\emph{every} program: $t_{\mathcal T} = t_{\mathcal D}$ for arbitrary
chains of neural symbols. By (i) no prediction changes; the price is one
$\mathrm{logsumexp}$ per head, exactly the operation whose deferral was
the point of \texttt{LogTens}.
\end{corollary}

This gives a dial with two labelled ends. Fully deferred (no
$\mathrm{nrm}$ anywhere): maximal numerical benefit, free exactly for
programs where nothing neural consumes a bound neural value, our MNIST
class. Fully normalized (a $\mathrm{nrm}$ after every head): safe for
every program; this is logLTN's $\log \circ \mathrm{softmax}$ composition
(cf.\ Sec.~\ref{sec:related}) adopted head by head, while keeping log
space as the native carrier. In between, the bound
$\sum_i \log(M_i/m_i)$ measures how far a partially deferred program can
drift. The theorem says where on the dial a given program shape must sit.

\item[T4: ablation, remove lazy evaluation (Rev.~2).]
Add a strict mode to the JAX backend (that is where all reported numbers
come from, l.529): build every \texttt{Bind} eagerly, full outer products,
no pruning. Run $N=1,2,4$ at batch $32$ and report step time and peak
memory next to the lazy default. Expected outcome: $N=4$ becomes impossible
or much slower, which is the concrete evidence Rev.~2 asks for. Paper: two
sentences and a small table after Table~\ref{tab:mnist}, or inside the
Longer numbers paragraph; then the laziness claim in
Sec.~\ref{sec:related} (l.603) can point at numbers.

\item[T5: vary the quantifier aggregation $p$ (Rev.~2).]
Quantifiers are aggregations
$\mathcal I(Q)_D \colon (D\to\mathcal M\,\Omega)\to\mathcal M\,\Omega$
(Sec.~\ref{sec:layer-categorical-logical}). Training currently averages
negative logs, which is a geometric mean (Sec.~\ref{sec:training}). Try
LTN's $p$-mean family instead (in log space: a generalized log mean exp)
for $p \in \{1,2,4,8\}$ on $N=1,2$, LTN style variant, $15$ seeds; report
accuracy and convergence. The question this answers: does LTN's $p$-mean
trick against getting stuck early help here too, or does exact marginal
likelihood make it unnecessary? Paper: a short Aggregation choices
paragraph in Sec.~\ref{sec:eval}.

\item[T6: truth space $\mathbb R_{\ge 0}$, log truth space, semirings (Rev.~3).]
Set up the tensor monad over any commutative semiring $S$: $\mathcal M_S X$
= finitely supported maps $X \to S$, bind = weighted pushforward (one short
proposition). Then use two instances: $(\mathbb R_{\ge 0},+,\cdot)$, the
truth space of unnormalized masses, and the log semiring
$(\mathbb R\cup\{-\infty\},\mathrm{logsumexp},+)$, the log truth space.
That makes $-\infty$ the real zero of log space, which answers Rev.~3's
missing log zero point (right now l.516 admits an $\epsilon$ floor stand
in). Also say that floating point numbers form an approximate semiring,
which separates the exact semantics from the finite precision
implementation. Places to edit: the Table~\ref{tab:monads-nesycat} row
(l.138 to 141) and the \texttt{LogTens} definition (l.423 to 425), both
currently over $\mathbb R$ without $-\infty$, plus the $\epsilon$ floor
remark (l.516). Note: T3 needs this, the unit law of the bridge only holds
exactly with the true $-\infty$. Also it is important to describe the logical operations on the log semiring, that come from the lift of Bool to T Bool.

\item[T7: related work, KLay (Rev.~2).]
Cite \citet{maeneKLAYACCELERATINGARITHMETIC2025}. The key already exists in
\texttt{daniel\_zotero.bib} (l.5253); fix its metadata (add the ICLR 2025
booktitle and the third author Zuidberg Dos Martires) instead of adding a
duplicate. Where: right after the sentence ending ``(Table~3).'' (l.602),
before DeepStochLog. What to say: KLay turns already compiled circuits
(d-DNNF or SDD) into layered scatter and gather operations that run well on
GPUs. Same goal, tensorized exact inference, but the compilation step
stays; with us there is nothing to compile, the bind does the
marginalization and laziness does the pruning.

\item[T8: abstract, tighten the results claim (Rev.~1).]
Rewrite one sentence (l.71). Old: ``while achieving nearly the accuracy of
DeepStochLog. However, unlike DeepStochLog, we stay in a uniform framework
that applies to many first-order NeSy approaches.'' New, roughly:
``outperform LTN and DeepProbLog in speed and accuracy; DeepStochLog and
A-NeSI stay slightly more accurate (Table~\ref{tab:addition}). Our
contribution is not top accuracy but generality: one monad parametric
framework that covers classical, fuzzy and probabilistic first order
semantics at competitive accuracy and speed.'' Keep the Giry sentence
(l.72).

\item[T9: speed claim across GPU generations (Rev.~1).]
Extend l.554 to 555 (``well under the $5.36$/$3.44$\,ms LTN reports on an
older V100 GPU'') with the missing control: on FP32 CNN training, published
benchmarks put the A100 at no more than $2$ to $3$ times a V100, so the
V100 is at least a third as fast. Our gap is $3.44/0.52 \approx 6.6$ and
$3.44/0.44 \approx 7.8$, bigger than any plausible hardware factor, so the
speed advantage is not a GPU artifact. One sentence plus a footnote with
the benchmark source. Note: the same CNN, same batch size, same epochs
statement is in the body (l.551 to 552), not in Table~\ref{tab:mnist}'s
caption.

\item[T10: remove the $0.99^{2N}$ Reference row (Rev.~1).]
Delete the Reference row of Table~\ref{tab:addition} (l.575 and 576, the
\verb|\midrule| plus the Reference line); the caption needs no change
(checked). Then rewrite the Longer numbers paragraph, which currently says
``(the \emph{Reference} row)'' at l.582 and would point at nothing. Put the
numbers inline, for example: ``a $99\%$ digit classifier caps sum accuracy
at $0.99^{2N}$, that is $92.27\%$ at $N{=}4$; \nesycattorch\ ($92.0$) lands
just below that cap, DeepStochLog ($92.7$) and A-NeSI ($92.6$) at it.''

\item[T11: DeepStochLog limits, A-NeSI factorization (Rev.~1).]
(i) After the DeepStochLog part (l.602 to 603, which already covers lost
mass on failing derivations and the changed aggregation), add once
Appendix B exists: ``Conversely, semantics outside the grammar derivation
reading, like the {\L}ukasiewicz instantiation of Appendix B, have no
DeepStochLog formulation at all.'' (ii) A-NeSI: the point is already in the
paper (l.605 to 606, ``requires additional factorization preparations,
which are example specific and therefore not generalizable''). Do not
repeat it, connect it: in the Longer numbers paragraph (l.584 to 585)
extend ``with DeepStochLog and A-NeSI a little ahead'' by ``at the price of
example specific factorization networks (A-NeSI) or a changed aggregation
semantics (DeepStochLog)''. (iii) Rev.~1 asks why anyone would prefer
generality: add two or three sentences to the Conclusion with a concrete
scenario, one specification used with probabilistic semantics for training
and with classical semantics for verification.

\item[T12: de-anonymise on acceptance (Rev.~1).]
Camera ready checklist: switch \verb|\documentclass[anon]| to
\verb|[final]| (l.1 and 2; the author block l.61 to 64 then shows up);
replace the three \texttt{anonymous.4open.science} links in the Conclusion
(l.618) with the real repositories; remove the \verb|\iffinalsubmission|
patch (l.51 to 56); add acknowledgements and funding. About $15$ minutes,
do it at acceptance.

\end{description}

\subsection*{Till}

\begin{description}

\item[T13: friendlier introduction of monads (Rev.~2).]
Put a short intuition bit before Definition~1 (Sec.~\ref{sec:monads}): a
monad is a type of computation plus a rule for chaining steps (the
programmable semicolon), then a tiny dice example,
\[
  \DO\{\, d_1 \leftarrow \mathit{die};\ d_2 \leftarrow \mathit{die};\
  \mathbf{return}\,(d_1 + d_2) \,\},
\]
which reads: the distribution of the sum of two dice. That is exactly our
MNIST axiom with dice instead of CNNs, so Sec.~\ref{sec:mnist} lands as a
payoff instead of a cold start. Also move the plain words reading of the
$\mathcal D$ bind (``first $x$ is drawn from $\rho$, then $y$ from
$f(x)$'', l.154) up next to the bind signature (l.114 to 117). About $12$
lines.

\item[T14: stress that MNIST $N{=}4$ is complex (Rev.~2).]
Extend the Longer numbers paragraph (l.580 to 585) and spell out the size:
$N{=}4$ means $8$ CNN calls and $10^{8}$ digit combinations per sample,
summed out exactly; DeepProbLog and LTN time out
(Table~\ref{tab:addition}); A-NeSI exists exactly because this size was
considered out of reach \citep{kriekenANeSI2023}. One or two sentences; can
be reused word for word in the rebuttal (T1).

\item[T15: complexity, memory, scalability (Rev.~2).]
Add a short Cost model paragraph at the end of Sec.~\ref{sec:training}
(after the rectangularity part, l.510 to 516): a leaf is a $[B,k]$ tensor,
so a chain of binds with support sizes $k_1,\dots,k_r$ costs
$O(B\,k_i\,k_{i+1})$ per step and $O(B\,\max_i k_i)$ memory. Work it out
for MNIST Add $N$ (digit supports $10$, sum support up to
$2\cdot 10^{N}{-}1$) and put the predicted scaling next to the measured
step times of Table~\ref{tab:mnist}. Derivation, three or four lines, in
the appendix; one pointer sentence in Sec.~\ref{sec:eval}.

\item[T16: fewer category words, more plain words (Rev.~2).]
Three small changes, no restructuring. (i) After each formal definition in
Sec.~\ref{sec:syntax-semantics} (domain signature, domain interpretation,
Definition~\ref{def:Semantics}) add one sentence starting ``In words:''
that says the same thing on the $\mathsf{digit}/{+}/{=}$ example, with a
pointer to Sec.~\ref{sec:mnist}. (ii) Move the technical footnotes (like
the relative monad footnote, l.351 to 353) to Appendix~\ref{app:cat}; the
Kleisli triple part is already there (l.637 to 642), only the footnotes are
left to move. (iii) Open Sec.~\ref{sec:syntax-semantics} with two plain
sentences: signatures declare symbols, interpretations give them meaning,
one induction computes truth. Optional: use the
deterministic/effectful/neural wording (already in the table, l.208 to
219) when Tarski/Kleisli/Nesy first appear.

\end{description}
\endgroup
\end{document}
